\documentclass[a4paper,11pt,oneside]{memoir}

%%%%%%%%%Paquetes 
\usepackage[spanish,es-tabla]{babel}
\selectlanguage{spanish}
\usepackage[utf8]{inputenc}
\usepackage{placeins}
\usepackage{float}

\RequirePackage{booktabs}
\RequirePackage[table]{xcolor}
\RequirePackage{xtab}
\RequirePackage{multirow}

%Paquetes adicionales
\usepackage{csquotes}
\usepackage{amsmath}
\usepackage{eurosym}
\usepackage[version=4]{mhchem}
\usepackage{soul}
\usepackage{textgreek}
\usepackage{adjustbox}





%Para establecer las imagenes exactamente debajo de los parrafos
\usepackage{float}

%Para codigo
\usepackage{listings}

% Multi-page tables using
\usepackage{longtable}
\usepackage{tabularx}

%Citas
\usepackage{cite}

%url
\usepackage{url}

%Para establecer las imagenes exactamente debajo de los parrafos
\usepackage{float}

% Ecuaciones
\usepackage{amsmath}

% Rutas de fichero / paquete
\newcommand{\ruta}[1]{{\sffamily #1}}

% Párrafos
\nonzeroparskip

%%%%%%%%%Bibliografia / Links
\usepackage[numbers,sort]{natbib}

% Links
\usepackage[colorlinks, backref=page]{hyperref}
\hypersetup{
	colorlinks,
	linkcolor={gray!40!black},
	citecolor={blue!50!black},
	urlcolor=blue,
	filecolor=magenta,
}

\renewcommand*{\backref}[1]{}
\renewcommand*{\backrefalt}[4]{
	\ifcase #1 (No citado.)
	\or (Citado en página~#2.)
	\else (Citado en páginas #2.)
	\fi%
}
\renewcommand*{\backrefsep}{, }
\renewcommand*{\backreftwosep}{ y~}
\renewcommand*{\backreflastsep}{ y~}
	

%%%%%%%%%Imagenes- Esquemas
\usepackage{graphicx}
\newcommand{\imagen}[2]{
	\begin{figure}[H]
		\centering
		\includegraphics[width=0.9\textwidth]{#1}
		\caption{#2}\label{fig:#1}
	\end{figure}
	\FloatBarrier
}

% Imagenes- Logos y otros
\newcommand{\imagenflotante}[2]{
	\begin{figure}[H]
		\centering
		\includegraphics[width=0.4\textwidth]{#1}
		\caption{#2}\label{fig:#1}
		\FloatBarrier
	\end{figure}
}

% El comando \figura nos permite insertar figuras comodamente, y utilizando
% siempre el mismo formato. Los parametros son:
% 1 -> Porcentaje del ancho de página que ocupará la figura (de 0 a 1)
% 2 --> Fichero de la imagen
% 3 --> Texto a pie de imagen
% 4 --> Etiqueta (label) para referencias
% 5 --> Opciones que queramos pasarle al \includegraphics
% 6 --> Opciones de posicionamiento a pasarle a \begin{figure}
\newcommand{\figuraConPosicion}[6]{%
  \setlength{\anchoFloat}{#1\textwidth}%
  \addtolength{\anchoFloat}{-4\fboxsep}%
  \setlength{\anchoFigura}{\anchoFloat}%
  \begin{figure}[#6]
    \begin{center}%
      \Ovalbox{%
        \begin{minipage}{\anchoFloat}%
          \begin{center}%
            \includegraphics[width=\anchoFigura,#5]{#2}%
            \caption{#3}%
            \label{#4}%
          \end{center}%
        \end{minipage}
      }%
    \end{center}%
  \end{figure}%
}

% Comando para incluir imágenes en formato apaisado (sin marco).
\newcommand{\figuraApaisadaSinMarco}[5]{%
  \begin{figure}%
    \begin{center}%
    \includegraphics[angle=90,height=#1\textheight,#5]{#2}%
    \caption{#3}%
    \label{#4}%
    \end{center}%
  \end{figure}%
}

%%%%%%%%%Tablas
\newcommand{\otoprule}{\midrule [\heavyrulewidth]}

% Nuevo comando para tablas pequeñas (menos de una página).
\newcommand{\tablaSmall}[5]{%
 \begin{table}
  \begin{center}
   \rowcolors {2}{gray!35}{}
   \begin{tabular}{#2}
    \toprule
    #4
    \otoprule
    #5
    \bottomrule
   \end{tabular}
   \caption{#1}
   \label{tabla:#3}
  \end{center}
 \end{table}
}

% Nuevo comando para tablas pequeñas (menos de una página).
\newcommand{\tablaSmallSinColores}[5]{%
 \begin{table}[H]
  \begin{center}
   \begin{tabular}{#2}
    \toprule
    #4
    \otoprule
    #5
    \bottomrule
   \end{tabular}
   \caption{#1}
   \label{tabla:#3}
  \end{center}
 \end{table}
}

\newcommand{\tablaApaisadaSmall}[5]{%
\begin{landscape}
  \begin{table}
   \begin{center}
    \rowcolors {2}{gray!35}{}
    \begin{tabular}{#2}
     \toprule
     #4
     \otoprule
     #5
     \bottomrule
    \end{tabular}
    \caption{#1}
    \label{tabla:#3}
   \end{center}
  \end{table}
\end{landscape}
}

% Nuevo comando para tablas grandes con cabecera y filas alternas coloreadas en gris.
\newcommand{\tabla}[6]{%
  \begin{center}
    \tablefirsthead{
      \toprule
      #5
      \otoprule
    }
    \tablehead{
      \multicolumn{#3}{l}{\small\sl continúa desde la página anterior}\\
      \toprule
      #5
      \otoprule
    }
    \tabletail{
      \hline
      \multicolumn{#3}{r}{\small\sl continúa en la página siguiente}\\
    }
    \tablelasttail{
      \hline
    }
    \bottomcaption{#1}
    \rowcolors {2}{gray!35}{}
    \begin{xtabular}{#2}
      #6
      \bottomrule
    \end{xtabular}
    \label{tabla:#4}
  \end{center}
}

% Nuevo comando para tablas grandes con cabecera.
\newcommand{\tablaSinColores}[6]{%
  \begin{center}
    \tablefirsthead{
      \toprule
      #5
      \otoprule
    }
    \tablehead{
      \multicolumn{#3}{l}{\small\sl continúa desde la página anterior}\\
      \toprule
      #5
      \otoprule
    }
    \tabletail{
      \hline
      \multicolumn{#3}{r}{\small\sl continúa en la página siguiente}\\
    }
    \tablelasttail{
      \hline
    }
    \bottomcaption{#1}
    \begin{xtabular}{#2}
      #6
      \bottomrule
    \end{xtabular}
    \label{tabla:#4}
  \end{center}
}

% Nuevo comando para tablas grandes sin cabecera.
\newcommand{\tablaSinCabecera}[5]{%
  \begin{center}
    \tablefirsthead{
      \toprule
    }
    \tablehead{
      \multicolumn{#3}{l}{\small\sl continúa desde la página anterior}\\
      \hline
    }
    \tabletail{
      \hline
      \multicolumn{#3}{r}{\small\sl continúa en la página siguiente}\\
    }
    \tablelasttail{
      \hline
    }
    \bottomcaption{#1}
  \begin{xtabular}{#2}
    #5
   \bottomrule
  \end{xtabular}
  \label{tabla:#4}
  \end{center}
}

\definecolor{cgoLight}{HTML}{EEEEEE}
\definecolor{cgoExtralight}{HTML}{FFFFFF}

% Nuevo comando para tablas grandes sin cabecera.
\newcommand{\tablaSinCabeceraConBandas}[5]{%
  \begin{center}
    \tablefirsthead{
      \toprule
    }
    \tablehead{
      \multicolumn{#3}{l}{\small\sl continúa desde la página anterior}\\
      \hline
    }
    \tabletail{
      \hline
      \multicolumn{#3}{r}{\small\sl continúa en la página siguiente}\\
    }
    \tablelasttail{
      \hline
    }
    \bottomcaption{#1}
    \rowcolors[]{1}{cgoExtralight}{cgoLight}

  \begin{xtabular}{#2}
    #5
   \bottomrule
  \end{xtabular}
  \label{tabla:#4}
  \end{center}
}

%Path imagenes
\graphicspath{ {./img/} }

% Capítulos
\chapterstyle{bianchi}
\newcommand{\capitulo}[2]{
	\setcounter{chapter}{#1}
	\setcounter{section}{0}
	\chapter*{#2}
	\addcontentsline{toc}{chapter}{#2}
	\markboth{#2}{#2}
}

% Apéndices
\renewcommand{\appendixname}{Apéndice}
\renewcommand*\cftappendixname{\appendixname}

\newcommand{\apendice}[1]{
	%\renewcommand{\thechapter}{A}
	\chapter{#1}
}

\renewcommand*\cftappendixname{\appendixname\ }

%%%%%%%%%Portada
\makeatletter
\usepackage{xcolor}
\newcommand{\tutor}[1]{\def\@tutor{#1}}
\newcommand{\course}[1]{\def\@course{#1}}
\definecolor{cpardoBox}{HTML}{E6E6FF}
\def\maketitle{
  \null
  \thispagestyle{empty}
  % Cabecera ----------------
\begin{center}%
	{\noindent\Huge Universidades de Burgos, León y Valladolid}\vspace{.2cm}%
	
	{\noindent\Large Máster universitario}\vspace{.5cm}%
	
	{\noindent\Huge \textbf{Inteligencia de Negocio y Big~Data en Entornos Seguros}}\vspace{.5cm}%
\end{center}%

\begin{center}%
	\includegraphics[height=3cm]{img/escudoUBU} \hspace{1cm}
	\includegraphics[height=3cm]{img/escudoUVA} \hspace{1cm}
	\includegraphics[height=3cm]{img/escudoULE} \vspace{1cm}%
\end{center}%

  \vfill
  % Título proyecto y escudo informática ----------------
  \colorbox{cpardoBox}{%
    \begin{minipage}{.9\textwidth}
      \vspace{.5cm}\Large
      \begin{center}
      \textbf{TFM del Máster Inteligencia de Negocio y Big Data en Entornos Seguros}\vspace{.6cm}\\
      \textbf{\LARGE\@title{}}
      \end{center}
      \vspace{.2cm}
    \end{minipage}

  }%
  \hfill
  \vfill
  % Datos de alumno, curso y tutores ------------------
  \begin{center}%
  {%
    \noindent\LARGE
    Presentado por \@author{}\\ 
     ---\\ \@date{}\\
    Tutor Académico:  \tutorA \\
    Tutor Empresa:  \tutorE \\
  }%
  \end{center}%
  \null
  \cleardoublepage
  }
\makeatother

%%% Declaración de variables 
\newcommand{\nombre}{Adrián Aguado García} %%% cambio de comando
\newcommand{\tutorA}{Carlos E.Vivaracho Pascual} %%% cambio de comando
\newcommand{\tutorE}{Fernando Cuenca Cabezas} %%% cambio de comando
\newcommand{\titulo}{Estudio analítico sobre la calidad del aire en la ciudad de Madrid} %%% cambio de comando


% Datos de portada
\title{\titulo}
\author{\nombre}
\tutor{\tutorE}
\tutor{\tutorA}
\date{\today}

\begin{document}
\maketitle

%Pagina en blanco
\newpage\null\thispagestyle{empty}
\newpage\null\thispagestyle{empty}

%%%%%%%%%%%%%%%%%%%%%%%%%%%%%%%%%%%%%%%%%%%%%%%%%%%%%%%%%%%%%%%%%%%%%%%%%%%%%%%%%%%%%%%%
\thispagestyle{empty}

\noindent
\begin{center}%
	{\noindent\Huge Universidades de Burgos, León y Valladolid}\vspace{.5cm}%
	
\begin{center}%
	\includegraphics[height=3cm]{img/escudoUBU} \hspace{1cm}
	\includegraphics[height=3cm]{img/escudoUVA} \hspace{1cm}
	\includegraphics[height=3cm]{img/escudoULE} \vspace{1cm}%
\end{center}%

	{\noindent\Large \textbf{Máster universitario en Inteligencia de Negocio y Big~Data en Entornos Seguros}}\vspace{.5cm}%
\end{center}%



\begin{center} \tutorA, profesor de la Universidad de Valladolid, del departamento de Informática, área de Ciencias de la Computación y la Inteligencia Artificial.

\noindent Expone:

\noindent Que el alumno D. \nombre, con DNI 78759314T, ha realizado el Trabajo final de Máster en Inteligencia de Negocio y Big Data en Entornos Seguros 
          titulado "\titulo". 

\noindent Y que dicho trabajo ha sido realizado por el alumno bajo la dirección del que suscribe, en virtud de lo cual se autoriza su presentación y defensa.
\end{center}%
\begin{center} %\large
En Burgos, {\large \today}
\end{center}

\vfill\vfill\vfill

% Author and supervisor
\begin{minipage}{0.45\textwidth}
\begin{flushleft} %\large
\begin{center}
Vº. Bº. del Tutor:\\[2cm]
D. nombre tutor\\
\textit{\tutorA}
\end{center}
\end{flushleft}
\end{minipage}
\hfill
\begin{minipage}{0.45\textwidth}
\begin{flushleft} %\large
\begin{center}
Vº. Bº. del co-tutor:\\[2cm]
D. nombre co-tutor\\
\textit{\tutorE}
\end{center}
\end{flushleft}
\end{minipage}
\hfill

\vfill

%Pagina en blanco
\newpage\null\thispagestyle{empty}\newpage




\frontmatter

% Abstract en castellano
\renewcommand*\abstractname{Resumen}
\begin{abstract}
La contaminación ambiental supone, y cada vez más, un problema para la salud humana en todo el mundo. Los núcleos urbanos grandes resultan ser una de las áreas con mayor contaminación. Más concretamente diversos gases, como el dióxido de nitrógeno \ce{NO2}, están haciendo que las administraciones se planteen diversas medidas para su reducción. En el presente documento se realiza un estudio basado en un modelo analítico para los datos de contaminación de la ciudad de Madrid.
\end{abstract}

\renewcommand*\abstractname{Descriptores}
\begin{abstract}
Big Data, Visualización de datos, Python, R, contaminación ambiental\ldots
\end{abstract}

\clearpage

% Abstract en inglés
\renewcommand*\abstractname{Abstract}
\begin{abstract}
Environmental pollution is a problem for human health worldwide. Large urban centres are one of the most polluted areas. More specifically, various gases, such as nitrogen dioxide \ce{NO2}, are causing administrations to consider various measures for their reduction. This document is a study based on an analytical model for pollution data from the city of Madrid.
\end{abstract}

\renewcommand*\abstractname{Keywords}
\begin{abstract}
Big Data, Data Visualization, Python, R, pollution \ldots
\end{abstract}

\clearpage

% Indices
\tableofcontents

\clearpage

\listoffigures

\clearpage

\listoftables
\clearpage

\mainmatter
\clearpage

\addcontentsline{toc}{part}{Memoria}
\part*{Memoria}

\capitulo{1}{Introducción}

La contaminación ambiental se ha incrementado, y sigue haciéndolo a un ritmo vertiginoso, de manera preocupante en los últimos años. Constituye uno de los problemas más serios a los que se enfrenta el ser humano. Hoy día ya no es una cuestión localizada en algunos lugares sino que el viento se ha encargado de convertirlo en un problema global. Los gases provenientes de automóviles, camiones, procesos industriales, sistemas de calefacción e incluso hasta el humo de los cigarrillos de miles de fumadores se juntan para contaminar el aire que consumimos a diario. 

Respirar aire limpio y sin riesgos para la salud debería ser un derecho de toda persona. Está demostrado que la contaminación atmosférica causa graves daños a la salud y al medio ambiente. Los niveles actuales de contaminación atmosférica provocan la muerte de entre cuatro y cinco millones de personas por exposición directa al aire contaminado en todo el mundo \cite{sharma_anirudh_anirudh_2018}. Si hablamos de España alrededor de 16.000 muertes prematuras en España \cite{informe_EeA_2019}.

Pero la contaminación es un tema realmente amplio, dado que la mayor parte de contaminantes son expulsados por procesos industriales y automóviles; y además éstos se concentran principalmente en grandes urbes, es precisamente ahí donde debemos realizar nuestros mayores esfuerzos. Casi la mitad de la población mundial vive actualmente en ciudades, y para el año 2050 se prevé que aumente a un 75\% \cite{bbc_ciudades}. Nuestro reto ahora es el intentar cambiar la manera en la que vivimos en aquellos puntos de mayor concentración de contaminación para evitar seguir contaminando de la manera que lo hacemos. Globalmente ya se están tomando medidas, diversos países proponen medidas para tratar de frenar o reducir la contaminación. 

En el caso concreto de la ciudad de Madrid, se trata de una de las dos ciudades españolas, junto con Barcelona, que está obligada a cumplir los niveles máximos de \ce{NO2} impuestos como medidas para reducir la contaminación dictados por la unión europea.

En las páginas siguientes vamos a realizar un estudio mediante el cuál analizaremos la calidad de aire de la ciudad de Madrid gracias al portal de datos abiertos y basándonos en los niveles de calidad del aire que recomienda la unión europea como normales.
\capitulo{2}{Objetivos del proyecto}

En este apartado se van a detallar los diferentes objetivos que se buscaban con la realización de este trabajo fin de máster.

Se parte de la hipótesis de que es posible obtener datos abiertos que permitan caracterizar  los niveles de contaminación de la ciudad de Madrid y que a partir de estos datos se puede crear un modelo para obtener, analizar y visualizar los datos 
de la calidad del aire.

\section{Objetivos}\label{objetivos-generales}

\begin{itemize}
\tightlist
\item Analizar, usando los datos abiertos que proporciona el Ayuntamiento de la ciudad de Madrid, la calidad del aire de la ciudad, poniéndola en relación con los objetivos planteados por la OMS.
\item Preparación de los datos (Selección, limpieza, transformación, preprocesado)
\item Visualización de datos que ayude a la respuesta de diversas preguntas analíticas 
\item Respuesta de diversas preguntas analíticas
\item Elección de unos contaminantes para analizar en detalle
\end{itemize}

\section{Objetivos personales}\label{objetivos-personales}
\begin{itemize}
\tightlist
\item
 Realizar un trabajo fin de máster más enfocado a la investigación. Durante trabajos académicos anteriores siempre había partido de una base mas técnica con unos objetivos finales puramente técnicos. En este caso la combinación de una parte técnica con una parte más de exploración o investigación hizo que el alumno se decantara por esta opción. La base del máster proporcionaba casi todo lo necesario para aplicar las técnicas aprendidas con un bonito reto a explorar. Además no descarto el realizar un doctorado en un futuro y creo que puede venir bien comenzar a tener esta visión más de investigación.
 \item
 Realización de un proyecto que involucre datos reales
\item
  Mejorar mis conocimientos de Python, R, visualización de datos y estadística. 
  \item Profundizar en el campo del \textit{data science}.
\end{itemize}

\capitulo{3}{Conceptos teóricos}

Una de las partes más importantes del proyecto, en cuestión de tiempo, ha sido la de familiarizarse con los conceptos científicos que requería el mismo. Además dada la numerosa cantidad de información sobre el tema en algunos momentos he tenido dudas sobre el punto de vista que darle al análisis dada la abrumadora cantidad de opciones de las que disponía. 

\section{Introducción}\label{teorico-introduccion}
Esta sección pretende servir como orientación antes de adentrarse en el trabajo en sí mismo. Se trata de un conglomerado de consideraciones de diferente índole que es necesario conocer para comprender este trabajo fin de máster. Esta subdividido en un preámbulo, algo más general y breve, sobre conceptos básicos de un proyecto \emph{Big Data}. Asimismo se abordan cuestiones como los tipos de datos, arquitecturas o smart cities. Se prosigue con diversos análisis de índole científica, legislativa y, por último, técnica. Todas estas secciones descritas siempre enfocadas desde el punto de vista del trabajo que se esta realizando.

\section{Big data}\label{teorico-big-data}
La información siempre ha sido, y será, el tesoro más valioso con el que puede contar una empresa. En la era de la sociedad de la información y en la que, al día, se generan 2,5 quintillones \footnote{Fuente IBM, Noticia ABC. \url{https://www.abc.es/tecnologia/redes/20140423/abci-trillones-byte-informacion-cada-201404222207.html}} de datos, resulta más que evidente que se trata de la herramienta empresarial y social más importante de la historia. 

Empresas de cualquier tamaño pueden y deben usar esta tecnología para cubrir prácticamente cualquier necesidad que tengan. Se trata de un nuevo prisma que modifica la visión de un negocio desde cualquier perspectiva. Sin embargo podemos afirmar categóricamente que se trata de un área de reciente implantación, que avanza a una velocidad vertiginosa, y que pocas empresas emplean con resultados altamente satisfactorios. Por todo ello no existen definiciones estrictamente formales del término. Una aproximación a una interpretación del término podría ser la siguiente:

\begin{displayquote}
	\emph{Big Data refers to a data set that is so large and/or complex that it cannot be perceived, acquired, managed, and processed by traditional Information Technology (IT) and software/hardware tools within a tolerable time} \cite{laval_introduction}
\end{displayquote} 

En otras palabra, \emph{Big data} es un concepto en auge que se puede describir como una gran masa de datos cuyo gran tamaño hace complejo su análisis con las técnicas más habituales. El problema actual reside en que, y vamos a más, los grandes volúmenes de datos han superado con creces a la capacidad de procesamiento de un simple host. La importancia del término mencionado nació para intentar solucionar esos problemas de almacenamiento e intentar interpretar, almacenar y analizar esas grandes cantidades de datos de la manera más adecuada y eficiente posible.

Cada día que nos levantamos si miramos a nuesrtro alrededor podemos ver que cualquier dispositivo de los que usamos a diario genera datos constantemente, todos ellos pueden ser analizados. De toda esta maraña de datos aproximadamente solo un 20\% es información estructurada, mientras que lo restante se trata de datos no estructurados. Éstos últimos requieren de un análisis un poco más profundo para almacenarlos y analizarlos. 

Por lo que vemos resulta bastante complejo el intentar dar una definición exacta, pero en lo que sí estamos de acuerdo es en que podríamos agrupar la acepción mediante tres dimensiones, conocidas como “las tres V”: Volumen, Variedad y Velocidad.

\begin{itemize}
	\tightlist
	\item
	\textbf{Velocidad}: cada vez más y más datos generados a ritmo trepidante.
	\item
	\textbf{Variedad}: referencia a los tipos de datos. Los datos pueden datos pueden ser, simplificando, de dos tipos: estructurados o no estructurados. La manera de abordar el proyecto cambia en función de los datos a analizar.
	\item
	\textbf{Volumen}: alusión a la capacidad de almacenaje de todos esos datos que estamos analizando
\end{itemize}

\imagen{3vbigdata}{\footnotesize{Las 3 v's en \emph{Big Data}. Fuente: \url{cadreresearch.org}}}

Éstas tres \textit{v's} engloban los principales atributos del \emph{Big Data}, aunque también es cierto que hoy en día y dada la creciente repercusión de las fake news y demás escándalos suelen incluir una cuarta, la \textbf{Veracidad}. Si nos encontráramos en un entorno \textit{business} también es posible incluir una quinta v, el \textbf{Valor}. Que no es más que el valor que tienen esos datos. 

\subsection{Datos}
Otra parte fundamental de el proceso que describimos es conocer que existen diferentes tipos de datos. Se pueden encontrar por lo tanto tres subgrupos de datos.
\begin{itemize}
	\tightlist
	\item	Datos \underline {estructurados}: es información ya procesada, filtrada y con un formato estructurado. Las tablas relacionales son un buen ejemplo.
	\item	Datos \underline {semi-estructurados}: es información que posee un formato definido y esta procesada pero no se encuentra estructurada. Un ejemplo serían los ficheros con información de un lenguaje de etiquetas como HTML.
	\item Datos \underline {no estructurados}: se trata de información sin procesar y que puede tener cualquier estructura. La información colgada en redes sociales serían un gran ejemplo de este tipo de datos.
\end{itemize}
La manera de trabajar hoy en día implica almacenar solamente datos de tipo estructurado o semiestructurado, obligando a pasar por un proceso de filtrado y transformación los datos no estructurados. Este proceso radica en un coste y en una pérdida inevitable de datos que cada vez es más difícil ignorar ya que cada vez se busca más precisión y además iría en contra de las \textit{v’s} nombradas anteriormente.

\subsection{Arquitecturas}

La arquitectura de un proyecto \emph{Big Data} está compuesta generalmente por varias capas, es posible encontrar fuentes que hablan de cuatro y también de cinco, aunque quizás es verdad que depende del proyecto al que se aplican. Particularmente creo que está mejor enfocado con cuatro capas que serían las siguientes: \emph{recolección de datos, almacenamiento, procesamiento de datos y visualización}. Esta arquitectura “por capas” no es nueva, sino que ya es algo generalizado en las soluciones de \textit{Business Intelligence} que existen hoy en día. Sin embargo, debido a las nuevas necesidades cada uno de estos pasos ha ido adaptándose y aportando nuevas tecnologías a la vez que abriendo nuevas oportunidades.

\begin{enumerate}
	\item \textbf{Recolección}: La primera etapa es dónde el sistema debe conectarse a las fuentes para extraer la información. Dependiendo de cómo se conecte a la fuente las herramientas pueden ser por lotes o por transmisión en tiempo real.
	\item \textbf{Almacenamiento}: En esta segunda fase se engloban los sistemas de ficheros y la base de datos.
	\item \textbf{Procesamiento}: Una vez se tienen los datos almacenados obtenidos del paso anterior, el siguiente paso en un sistema Big Data es explotar la
	información para llegar a los resultados deseados
	
	\item \textbf{Visualización}: Este es el apartado o paso es que menos ha cambiado si lo comparamos con otro tipo de arquitecturas. Es el paso en el que se muestran los datos de manera visual.
\end{enumerate}

\imagen{arquitecturaBData}{\footnotesize{Arq. tradicional vs Arq. \emph{Big Data}. Fuente: \url{https://www.slideshare.net/bd4s/big-data-introduccion}}}


\subsection{\emph{Smart cities}}

Aunque es práctica invisible para nosotros, cantidades ingentes de información se almacenan y circulan alrededor de nuestro mundo cada día generado mayormente por las empresas, administraciones o los propios ciudadanos en una explosión de datos a la espera de ser procesados y aprovechados con fines diversos. Desde hace unos años, la transformación digital que estamos sufriendo todos supone un cambio de paradigma en nuestras vidas. Existen multitud de actividades, por no decir casi todas, dónde la influencia del análisis de datos nos proporciona predicciones realmente interesantes. Si este tipo de actividades se realizan dentro de las ciudades, es decir, se aplica a la habitabildiad y al diseño de las mismas, entonces estaremos hablando de lo que los expertos llaman \emph{Smart Data} \cite{foote_big_2018}.

Las ciudades ofrecen distintos servicios a los ciudadanos por lo cual necesitan recopilar y almacenar una gran cantidad de datos. Otros organismos y administraciones públicas, financiados por la ciudadanía, generan a su vez información geográfica, cartográfica, meteorológica, médica etc. Los partidarios del \emph{Open Data}, la filosofía de datos abiertos, defienden que esta información debería ser accesible y reutilizable para el público en general, sin exigencia de ningun permiso específico. El por qué parece bastante lógico ya que consideran que restringir su acceso va contra el bien común debido a que se trata de información que pertenece a la sociedad, o ha sido financiada por ella. Por lo tanto, intentan promover que ese tipo de datos esten accesibles para todo el mundo que así lo desee. 

Hablamos de \textit{Smart Cities} cuando las ciudades tratan de innovar para mejorar la calidad de vida de sus habitantes. Se trata de un término cada vez má común y hace referencia a sistemas de infraestructura, conectividad, \emph{IoT} (Internet of things), y, por supuesto, datos abiertos que conviven interconectados entre sí en un núcleo de población. Hablamos también de tecnologías\textit{ Big Data}, \textit{machine learning} y de extracción, transformación, normalización, procesamiento y exposición de datos, que nos permitirán sacar partido a toda esa información, detectar problemas, extraer patrones, analizar comportamientos, plantear hipótesis u optimizar modelos. En definitiva, usar los datos disponibles para optimizar los servicios ofrecidos a los ciudadanos. Este término de \textbf{ciudad inteligente} se puede disgregar en dos partes:

\begin{itemize}
	\item Por un lado hace referencia al creciente concepto de que las grandes urbes están compuestas de una elevada digitalización. Prácticamente desde el proceso o tarea más básico se pretende digitalizar para proporcionar a la administración la vía para integrar y analizar la información producida por los ciudadanos y de esta manera realizar los procesos de dicha ciudad de una manera eficiente en base a esos datos producidos. Por ejemplo, el volumen de datos producidos por estos procesos puede ser utilizado para modelar y predecir el tráfico de las calles principales en una ciudad.
	\item Por otro lado, una \textit{smart city} es aquella cuya economía y gobierno está impulsado por la innovación, la creatividad y el emprendimiento. Quizás está parte sea la menos visible pero también resulta importante ya que las ciudades inteligentes no se basan todo en modelos tecnológicos sino que para evolucionar deben también cambiar la manera en la que funciona la administración. Por supuesto, las tecnologías de la información resultar ser de vital importancia, vistas como la plataforma para movilizar y llevar a cabo las ideas e innovaciones que se pretenden seguir para optimizar procesos.
\end{itemize}

Si queremos conocer algunas de estas ciudades nada mejor que visitar el informe anual de la escuela de negocios IESE, llamado \href{https://citiesinmotion.iese.edu/indicecim/}{IESE Cities in Motion}. Un gran ejemplo de ese tipo de ciudades es Nueva york.

\subsection{\emph{Big Data} y ciudades}
Hasta ahora nunca habíamos obtenido y recogido tal cantidad de datos en las ciudades, y lo que es quizás más importante, tenemos la capacidad para procesarlos y entenderlos. En las ciudades, aunque quizás jamás nos hayamos parado a pensarlo, es usual la producción de conjuntos de datos enormes. Además esto no resulta específico de hoy en día, sino desde hace tiempo se realizan censos, encuestas, entrevistas, baremos... El ser humano y sus acciones generamos datos pero las administraciones también. De la misma manera las empresas generan y almacenan cantidad de datos provenientes de sus operaciones internas. Sin embargo estamos hablando de que todos estos datos suelen ser muy específicos y generales por lo que resulta ser complejos a la hora de su análisis. Dependiendo de la manera de cómo se generan estos datos, es decir del origen que tienen dentro de las ciudades, podríamos hacer una clasificación de la siguiente manera:


\begin{itemize}
	\item Dirigido: aquí englobamos datos obtenidos median formas tradicionales de observación. En este tipo siempre suele existir una persona dentro de la iteracción que es la que normalmente almacena la información en una determinada base de datos. Algunos sería los controles de inmigración en los aeropuertos o los registros médico electrónicos.
	\item Automatizado: los datos obtenidos de manera automática son generados diariamente por sistemas basados en tecnología cuya función es captar datos de un determinado aspecto. Por ejemplo temperaturas medidas por sensores, transacciones de pago en tiendas o registros de comunicaciones.
	\item Voluntario: estos quizás sean los datos más evidentes y la vez no tanto. Son los datos que los usuario colgamos de manera voluntaria en nuestras redes sociales. Se puede extraer información muy valiosa de este tipo de datos. Algunos ejemplos son las publicaciones de Facebook, opiniones en twitter, búsquerda de Google o información sobre links. 
\end{itemize}


 Como ya hemos nombrado, la población humana presenta una clara tendencia a concentrarse en los núcleos urbanos \footnote{Fuente: Banco Mundial. \url{https://datos.bancomundial.org/indicador/sp.urb.totl.in.zs}}. Esta cantidad de personas hace que sea necesario el actuar en base a estos datos para intentar mejorar la calidad de vida de las personas que viven en esas ciudades. A mayor número de personas también crece el número de problemas y por lo tanto debemos utilizar la tecnología, en este caso la que hace referencia al \textit{Big Data}, para intentar reducir al mínimo u optimizar esos problemas. 
 
 Hoy en día las administraciones de grandes ciudades, en España tenemos ejemplos como Barcelona, Sevilla o Madrid, apuestan por un análisis en tiempo real de diferentes métricas relativas a las ciudades. Uno de los ejemplos más útiles es el del transporte público en tiempo real proporcionando a los ciudadanos implicacioens directas ya que pueden observar el estado de su próximo autobús o metro en una parada determinada y a una hora concreta. 
 
 En cualquier proyecto que engloben datos es necesario el obtener esos datos para posteriormente almacenarlos, procesarlos y analizarlos con el fin de obtener conocimiento en forma de información que resulte relevante para cadad ciudad. El primer próposito es el de conocer la finalidad con la que queremos recolectar unos datos determinados para así proveer a la ciudad de un nuevo sistema. Existen numerosas áreas en las que la ciudades pueden aprovechar sus datos, como hemos nombrado el aumento del número de personas genera problemas que que repercuten directamente en la calidad de vida de los habitantes de las ciudades, tales como la polución, el acceso de las personas a recursos y servicios, la congestión del tráfico, la seguridad y la movilidad de personas y mercancías, entre otros. A continuación se nombran algunas de las más relevantes. En la lista se incluyen tanto sitemas que se están desarrollando actualmente como propuestas que se pueden realizar, o se están realizando pero no he encontrado ejemplos, en diversas ciudades gracias al análisis de los datos que en ellas se generan.
 
 \begin{itemize}
 	\item \underline{Sistemas de transporte}: los problemas de movilidad, con el caso de transporte público o la congestión de tráfico, es posible paliarlos por medio de sistemas inteligentes de regulación de tráfico. Entre los medios utilizados para recabar datos de los agentes involucrados en el tráfico en las ciudades y de las infraestructuras viarias, destacan las videocámaras y diferentes tipos de dispositivos como sensores piezoeléctricos, magnéticos, inductivos, de infrarrojos, de microondas, radar, acústicos, e incluso sistemas que detectan dispositivos móviles e inalámbricos activos en el interior de los vehículos.
 	\item \underline{Seguridad}: sin duda es una de las ventajas más directas ya que través de cámaras de videovigilancia o sensores es posible controlar a los delincuentes de una manera más eficaz. Londres por ejemplo cuenta con más de 40.000 cámaras de vigiliancia aunque hay quien critica que quizás no sea la metodología más adecuada \cite{london_cameras}
 	\item \underline{Gestión de residuos}: en diversas ciudades ya se utilizan sistemas inteligente para optimizar la gestión de residuos tanto para la parte de clasificación como para la parte de recogida. \cite{big_data_ciudades_inteligentes}
 	\item \underline{Sancione}s: una de las aplicaciones más rentables para las adminitraciones en sin duda la del control del cumplimiento de las normas de ciculación. Algunos ejemplos que actualmente se pueden realizar de manera automática son, por ejemplo, invadir carriles reservados, pisar una línea continua, incumplir una señal de stop o saltarse un semáforo en rojo. Cámaras con reconocimiento facial y de matrículas permiten identificar el vehículo y el conductor infractor.
 	\item \underline{Sistemas medioambientes}: todo lo referente al medio ambiente se está utilizando desde hace años y se va a continuar utilizando para tratar de optimizarlo al máximo. Algunos de estos sistemas utilizan tecnología ubicua \cite{tec_ubicua}, capaz de suministrar información sobre diversas magnitudes de interés.
 	\item \underline{Energía eléctrica}: las empresas privadas de energía ya llevan años recopilando datos gracias a los contadores aunque desde hace unos años ya han empezado a surgir los contadores inteligentes que empiezan a realizar estimaciones de lo eficientes que somos a la hora de gastar energía.
 	\item \underline{Sistema de abastecimiento de agua}: el derroche de agua es un problema en muchas ciudades y con sistemas basados en datos extraídos a partir de sensrores de presión sería posible controlar la cantidad y calidad de agua en tiempo real.
 	\item \underline{Eficacia de servicios púbico}s: se me ocurre también la eficacia en actuaciones de cuerpos de seguridad, ambulancias o bomberos a través de la correlación de toda la información procedente de diversos sistemas y accesible a través de herramientas en tiempo real.
 \end{itemize}
 
Como vemos hay multitud de áreas de mejora para las ciudades inteligentes. Solo es cuestión de que las administraciones se planteen que es lo que se quiere optimizar para empezar a recabar datos y buscar la manera más óptima de utilizarlos en beneficio de la sociedad. 

En el caso de este trabajo nos centraremos en la parte de sistemas medioambientales o emisiones nocivas. Mediante sensores ubicados por toda la geografía de una ciudad se recogen de manera periódica diversos datos sobre diferentes gases que resultan ser nocivos para la salud. Asimismo también existen medidores de temperatura, precipitaciones u otro tipo de agentes medioambientales. En el apartado siguiente haremos una aproximación a la parte científica, de la parte de gases nocivos, necesaria para comprender mejor lo que miden este tipo de sensores para poder después abordarlo de manera práctica. El trabajo se centra en la ciudad de Madrid pero el análisis científico es válido para cualquier ciudad del mundo. Si bien es cierto que dependiendo de la zona en la que nos encontremos se tienen más en consideración un tipo u otro de compuestos químicos debido a la regulación existente en cada zona del mundo.


\section{Análisis científico}\label{teorico-cientifico}

Se define la contaminación atmosférica como \textit{“la presencia en la atmósfera de materias, sustancias o formas de energía que impliquen molestia grave, riesgo o daño para la seguridad o la salud de las personas, el medio ambiente y demás bienes de cualquier naturaleza}”.

La atmósfera es un bien común indispensable para la vida. Dada su condición esencial para la vida humana y por los daños que su contaminación o deterioro puede causar, la calidad del aire y la protección de la atmósfera ha sido, desde hace décadas, una prioridad política ambiental a lo largo de todo el mundo. Podemos afirmar que la contaminación atmosférica es consecuencia directa de las emisiones al aire de los gases y material particulado derivados de la actividad humana (social y económica) y de fuentes naturales \cite{informe_2018}.

Si hablamos del mundo, China lidera el ranking de países emisores de dióxido de carbono (\ce{CO2}), seguido de EE. UU. y de los estados miembros de la UE. El potente gas de efecto invernadero es una de las sustancias que más contribuyen al calentamiento global y al cambio climático, pero no es la única. Muchos otros compuestos, conocidos como contaminantes del clima, influyen en la cantidad de energía solar que retiene la Tierra... y provocan no pocos problemas de salud.

\imagen{graficaEmisiones}{\footnotesize{Emisiones de \ce{CO2} mundiales. Fuente: Infografía por \href{https://twitter.com/countcarbon/}{@countcarbon}}}

La gráfica que podemos ver en el párrafo inmediatamente superior  muestra qué países contaminan y cuánto: en vertical están las toneladas de \ce{CO2} por persona y en horizontal la población de cada país, de modo que el rectángulo simboliza las emisiones totales. China se lleva la palma, aunque sus emisiones son la mitad por persona que las de Estados Unidos. Hasta Alemania o Japón están peor (pero tienen menos personas y emisiones totales). 

En cuanto a europa, las últimas estimaciones globales de la AEMA (\href{https://www.eea.europa.eu/}{Agencia Europea de Medio Ambiente}) y la OMS (\href{https://www.who.int/}{Organización Mundial de la Salud}) sobre la repercusión sanitaria de la contaminación atmosférica son muy preocupantes. Elevan en el año 2015 hasta medio millón las muertes prematuras en los países europeos por la mala calidad del aire, 422.000 por exposición a partículas inferiores a 2,5 micras de diámetro (\ce{PM2,5}), 79.000 por exposición a dióxido de nitrógeno (\ce{NO2}) y 17.700 por exposición a ozono troposférico. En España, las víctimas de la contaminación serían ya más de 30.000 al año, 27.900 por partículas \ce{PM2,5}, 8.900 por \ce{NO2} y 1.800 por ozono, lo que supone duplicar los 16.000 fallecimientos prematuros anuales que se estimaban hace apenas una década \cite{informe_ecologistas}.

El coste económico de la mortalidad prematura y de la pérdida de días de trabajo por la contaminación del aire ambiente y en el interior de las viviendas ha sido cuantificado por el Banco Mundial en 38.000 millones de euros en 2013, equivalentes al 3,5 por ciento del Producto Interior Bruto (PIB) español, sin considerar los daños provocados a los cultivos, los ecosistemas naturales u otros bienes de cualquier naturaleza  \cite{informe_ecologistas}.



 \subsection{¿Cómo se mide la calidad del aire?}
 
 La calidad del aire viene determinada principalmente por la distribución geográfica de las fuentes de emisión de contaminantes y las cantidades de contaminantes que se emiten. Pero ¿qué es un contaminante? Según la directiva europea \textit{2008/50/CE} (CE, 2008), un contaminante es “toda sustancia presente en el aire ambiente que pueda tener efectos nocivos para la salud humana y el medio ambiente en su conjunto”. 
 
 Cualquier método que permita medir, calcular, predecir o estimar las emisiones, los niveles o los efectos de la contaminación atmosférica resulta válido para llevar a cabo la evaluación de la calidad del aire. Sin embargo dependiendo de los tipos de contaminantes en los que se quiera hacer foco existen métodos diferentes. Esto es debido a qué no todas las metodologías aportan la misma precisióm, se tienen en cuenta diversos factores.
 
   A finales de 2017, la Agencia Europea de Medio Ambiente (\href{https://www.eea.europa.eu/}{AEMA}) y la Comisión Europea pusieron en marcha un nuevo Índice europeo de Calidad del Aire (\href{https://www.eltiempo.es/calidad-aire}{ICA}), que permite a los usuarios comprobar su calidad en las más de 2.000 estaciones de medición repartidas por toda Europa. Este índice, que proporciona información actualizada sobre la calidad del aire en los 33 países miembros de la AEMA, incluye perfiles nacionales ya que las administraciones públicas locales tienen que adaptar sus medidas para considerar factores como demografía, infraestructuras de transporte, etc. Estas mediciones vigilan los estándares de calidad del aire y controlan los niveles de ozono (\ce{O3}), de dióxido de nitrógeno (\ce{NO2}), de dióxido de carbono (\ce{CO2}), de dióxido de azufre (\ce{SO2})... y toda la contaminación generada por partículas que pueden representar serios riesgos para la salud.
      
  La agencia medioambiente \href{https://www.epa.gov}{EPA} es la responsable del índice en EE. UU. mientras que el Centro Nacional para la Monitorización del Medio ambiente en China (CNMMC) es el organismo responsable de compilar, analizar, agregar y publicar los datos de los diferentes indicadores del aire en el país asiático.
  
   \imagen{ICAtabla}{\footnotesize{Tabla ICA. Fuente: \href{https://www.elespanol.com/omicrono/tecnologia/20171026/mide-contaminacion-aire-peligroso-podemos-hacer/257225717_0.html}{https://www.elespanol.com/omicrono}}
   \label{fig:tablaICA}}
   
   Esta contaminación suele darse  cuando las condiciones atmosféricas son favorablemente agradables (cuando hay Sol y el cielo está despejado). En estos casos, el suelo se calienta durante el día y se enfría durante la noche, de modo que el aire se estanca y no se regenera. El diseño de la ciudad también puede hacer aumentar o disminuir la polución, haciendo que el aire se estanque o se regenere en mayor o menor medida.
   
   El método técnico para saber la cantidad de contaminación que hay es mediante estaciones meteorológicas, también conocidas como estaciones de seguimiento de contaminación o estaciones remotas de medición de la calidad del aire.
   
   Estas estaciones miden la concentración de distintos agentes contaminantes en el aire. Además, se emplea la geoestadística, una ciencia no exacta, para predecir los niveles de contaminación por óxido de nitrógeno, una de las sustancias más perjudiciales para los seres vivos. Concretamente el dióxido de nitrógeno es producido por la combustión de los distintos combustibles que usan los vehículos mecánicos, por lo que es uno de los gases con mayor concentración.
      
   \begin{table}[H]
   	\begin{center}
   		\begin{tabular}{|l|l|}
   			\hline
   			Contaminante a analizar & Descripción medidores \\
   			\hline \hline
   			Dióxido de azufre (\ce{SO2}) & \begin{tabular}[c]{@{}l@{}}Este sistema se basa en la fluorescencia \\ producida por las moléculas del dióxido de azufre.\end{tabular}  \\ \hline
   			Monóxido de carbono (\ce{CO}) &  \begin{tabular}[c]{@{}l@{}}Este sistema se basa en la radiación infrarroja \\  que absorben las moléculas de monóxido de carbono. \end{tabular}	\\ \hline
   			Ozono (\ce{O3})  & Basado en la radiación ultravioleta. \\ \hline
   			Óxido de nitrógeno (\ce{NO}) &  \begin{tabular}[c]{@{}l@{}}Se basan en la energía que libera \\  la unión de óxido de nitrógeno con ozono. \end{tabular}	\\ \hline
   			Partículas en suspensión (\ce{PM2,5}) & \begin{tabular}[c]{@{}l@{}}Este tipo de dispositivos analizan partículas \\   en suspensión cuyo diámetro sea mayor a 2,5 \textmugreek m. \end{tabular}	\\ \hline
   			Hidrocarburos (\ce{CH})  & \begin{tabular}[c]{@{}l@{}} 	Estos dispositivos detectan las partículas \\ que libera la combustión del hidrógeno. \end{tabular}	\\ \hline
   		\end{tabular}
   		\caption{Tipos de medidores y sus contaminantes principales \cite{_sistema_vigilancia_madrid}}
   	\end{center}
   \end{table}   

 
  \subsection{Agentes contaminantes principales}
 
 Hablando en sentido general, los contaminantes no siempre son producidos por el hombre sino también por algunos fenómenos naturales, por ejemplo las erupciones volcánicas. En este breve apartado se va a incluir una lista con los agentes contaminantes principales \cite{_sistema_vigilancia_madrid}. Según su procedencia los podemos agrupar como sigue a continuación.
 
 \begin{itemize}
 	\item \underline{Primarios}: proceden de las fuentes de emisión
 	\begin{itemize}
 		\item \textit{Gaseosos}
 			\begin{itemize}
 				\item 	Dióxido de azufre (\ce{SO2})
 				\item 	Monóxido de carbono (\ce{CO})
 				\item 	Óxidos de nitrógeno (\ce{NOx})
 				\item 	Hidrocarburos (\ce{HC})
 				\item 	Dióxido de carbono (\ce{CO2})
 		    \end{itemize}
 		\item \textit{No gaseosos}
 			\begin{itemize}
 			\item Partículas: de índole  natural, en su mayoría, su procedencia es muy variada
 		   \end{itemize}
 	   	\item \textit{Otros}
 	   \begin{itemize}
 	   	\item 	  Restos de combustión de fuel, gas-oil o alquitranes
 	   	\item 	Erupciones volcánicas
 	   	\item 	Incendios
 	   	\item 	Intrusiones de material particulado
 	   	\item 	 Incineraciones no depuradas de basuras
 	   	\begin{itemize}
 	   		\item Metales pesados: plomo, cadmio, mercurio\ldots etc.
 	   	\end{itemize}
 	   \end{itemize}
 	\end{itemize}
 	\item \underline{Secundarios}: originados en la atmósfera como consecuencia de reacciones químicas que transforman los contaminantes primarios
 	\begin{itemize}
 \item	Ozono (\ce{O3})
 \item	Trióxido de azufre (\ce{SO3})
 \item	Ácido sulfúrico (\ce{H2SO4})
 \item	Dióxido de nitrógeno (\ce{NO2})
 \item	Ácido nítrico (\ce{HNO3})
 	\end{itemize}
 \end{itemize}

\subsection{La calidad del aire en España}

Todo este apartado esta basado en los diferentes informes que se realizan de manera anual por el ministerio de Agricultura, Alimentación y Medio Ambiente \cite{informe_2018}. Las siguientes líneas, que describen el marco teórico de la situación de España en cuanto a calidad de aire, se apoyan en el mencionado informe y también en el informe anual que elabora la plataforma ecologistas en acción \cite{informe_ecologistas}.

 En España, se divide el territorio en zonas o aglomeraciones, en función de la densidad de población, y se evalúa la calidad del aire para los contaminantes dióxido de azufre (\ce{SO2}), dióxido de nitrógeno y óxidos de nitrógeno (\ce{NO2}, \ce{NOX}), partículas (\ce{PM10} y \ce{PM2,5}), plomo (\ce{Pb}), benceno (\ce{C6H6}), monóxido de carbono (\ce{CO}), arsénico (\ce{As}), cadmio (\ce{Cd}), níquel (\ce{Ni)}, benzo(a)pireno (\ce{B(a)P}) y ozono (\ce{O3}), que son los contaminantes con valores legislados para protección de la salud.
 
 Para garantizar que se abarca la totalidad de la superficie nacional, las comunidades autónomas son las encargadas de dividir su territorio en zonas de calidad del aire homogéneas para la gestión y la evaluación (mediante mediciones, modelización u otras técnicas). Para ello se determinan unos métodos y criterios comunes de evaluación. También hay que cumplir con el requisito imprescindible de informar a la población y a las organizaciones interesadas. Por lo tanto cada comunidad es la que se encarga de las mediciones particulares de cada territorio. Después, al terminar el año. el gobierno se encarga de realizar el mencionado informe. El resultado de esta evaluación anual se presenta en uncuestionario técnico para su envío a la Comisión Europea y en otros informes más claros y comprensibles dirigidos a la población.
 
     \imagen{mapaTotalEstaciones}{\footnotesize{Mapa estaciones calidad del aire España.} Fuente: \href{https://www.mapama.gob.es/ide/metadatos/srv/spa/metadata.show?uuid=6c181182-6d75-4380-9783-d87a115d422a}{Metadatos gobierno España}}
 
 
 Por motivos obvios se va a escoger como referencia el último de los informes, que es el correspondiente al periodo anual del año anterior al que nos encontramos. Los contaminantes más problemáticos en el Estado español \textbf{durante 2018} han sido las partículas en suspensión (\ce{PM10} y \ce{PM2,5}), el dióxido de nitrógeno (\ce{NO2}), el ozono troposférico (\ce{O3}) y el dióxido de azufre (\ce{SO2}). Para el cálculo del porcentaje de población española  que respiró aire contaminado y de la superficie expuesta a niveles de contaminación que dañan la vegetación se han tenido en cuenta estos contaminantes, si bien se ha recopilado y evaluado asimismo la información disponible sobre otros contaminantes
 regulados legalmente como el monóxido de carbono (CO), el benceno, los hidrocarburos
 aromáticos policíclicos (HAP) y los metales pesados (arsénico, cadmio, níquel y plomo).
 
    \imagen{resumen2018contaminantes}{\footnotesize{Resumen calidad del aire por contaminante.} Fuente: Informe Calidad Aire 2018 \cite{informe_2018}}
 
Como conclusión general de ambos informes, y en referente a la salud de la población, se destacan dos perspectivas:
\begin{enumerate}
	\item La población que respiró aire contaminado en el Estado español, según los valores límite y objetivo establecidos para los contaminantes principales citados por la \textit{Directiva 2008/50/CE} y el Real Decreto \textit{102/2011}, alcanzó los 14,9 millones de personas, es decir un 31,8\% de toda la población. En otras palabras, \textbf{uno de cada tres españoles respiró un aire que incumple los estándares legales vigentes}. Esta situación supone no obstante un descenso de 2,6 millones de afectados respecto a 2017, y la cifra más baja desde el año 2011.
	\item Si se tienen en cuenta los valores recomendados por la Organización Mundial de la Salud (OMS), más estrictos que los límites legales (y más acordes con una adecuada protección	de la salud), \textbf{la población que respiró aire contaminado se incrementa hasta los 45,2 millones de personas}. Es decir, un 96,8\% de la población. En otras palabras, la práctica totalidad de los españoles respiró un aire con niveles de contaminación superiores a los recomendados por la OMS. Esta situación supone un modesto descenso de 0,6 millones de afectados respecto a 2017, y se mantiene por encima de la incidencia en la década, salvo el año 2015.
\end{enumerate}

\subsubsection{Metodologías de evaluación}

Independientemente de la comunidad autónoma o territorio  la evaluación de la calidad del aire debe efectuarse con un enfoque común basado en criterios de evaluación también comunes. Dicha evaluación debe tener en cuenta el tamaño de las poblaciones y los ecosistemas expuestos a la contaminación atmosférica, lo que lleva a clasificar el territorio nacional en zonas o aglomeraciones en función de la densidad de población \cite{casado_reinaldos_big_2018}.

\begin{itemize}
	\item Las zonas son porciones de territorio delimitadas por la Administración competente
	en cada caso utilizada para evaluación y gestión de la calidad del aire.
	\item Las aglomeraciones se definen como conurbaciones de población superiores a 250.000	habitantes o bien, cuando la población sea igual o inferior a 250.000 habitantes, con	una densidad de población por km2 que determine la Administración competente y
	justifique que se evalúe y controle la calidad del aire ambiente.
\end{itemize}

En las zonas y aglomeraciones así definidas se evalúa la calidad del aire para los contaminantes dióxido de azufre (\ce{SO2}), dióxido de nitrógeno y óxidos de nitrógeno (\ce{NO2}, \ce{NOx}),partículas (\ce{PM10} y \ce{PM2,5}), plomo (\ce{Pb}), benceno (\ce{C6H6}), monóxido de carbono (\ce{CO}), arsénico (\ce{As}), cadmio (\ce{Cd}), níquel (\ce{Ni}), benzo(a)pireno (\ce{B(a)P}) y ozono (\ce{O3}). La unidad utilizada en la medida de todos estos contaminantes es \textmugreek g/m3, microgramos por metro cúbico.

Dicha evaluación se efectúa considerando diversos objetivos de calidad del aire. 
Se distingue entre:
\begin{itemize}
	\item \underline{Valor límite}: Objetivo para la protección de la salud, definidos para \ce{SO2}, \ce{NO2},
			partículas \ce{PM10} y \ce{PM2,5}, plomo, benceno y \ce{CO}.
	\item \underline{Valor objetivo} (objetivos a largo plazo): Objetivos para la protección de la salud,	definidos para partículas \ce{PM2,5}, arsénico (\ce{As}), cadmio (\ce{Cd}), níquel (\ce{Ni}), \ce{B(a)P} y ozono  (\ce{O3}).
	\item	\underline{Nivel crítico}: Objetivos para la protección de la vegetación, definidos para \ce{SO2} y
	\ce{NOx}.
\end{itemize}

Se entiende por valor límite aquel fijado basándose en conocimientos científicos, con el
fin de evitar, prevenir o reducir los efectos nocivos para la salud humana, para el medio
ambiente en su conjunto y demás bienes de cualquier naturaleza que debe alcanzarse en
un período determinado y no superarse una vez alcanzado

Si nos centramos en el temas únicamente referidos a la salud existen unos \textbf{valores límite} objetivo que se recogen en la siguiente tabla: 
 
 %%%TABLE
\begin{table}[H]
	\begin{tabular}{|l|l|l|}
		\hline
		\multicolumn{1}{|c|}{\textbf{Contaminante}} & \multicolumn{1}{c|}{\textbf{Período de promedio}} & \multicolumn{1}{c|}{\textbf{Valor  límite}}                               \\ \hline
		\multirow{2}{*}{\ce{SO2}}                  & Horario                                           & 250 \textmugreek g / m3, (máx. 24 sup. al ano)                              \\ \cline{2-3} 
		& Diario                                            & \begin{tabular}[c]{@{}l@{}}50 \textmugreek g/m3\\ (máx. 3 sup. al año)\end{tabular}   \\ \hline
		\multirow{2}{*}{\ce{NO2}}                  & Horario                                           & \begin{tabular}[c]{@{}l@{}}200 \textmugreek g/m3\\ (máx. 18 sup. al año)\end{tabular} \\ \cline{2-3} 
		& Diario                                            & 40 \textmugreek g/m3                                                                  \\ \hline
		\multirow{2}{*}{\ce{PM10}}                 & Diario                                            & \begin{tabular}[c]{@{}l@{}}50 \textmugreek g/m3\\ (máx. 35 sup. al año)\end{tabular}  \\ \cline{2-3} 
		& Anual                                             & 40 \textmugreek g/m3                                                                  \\ \hline
		\ce{Pb}                                    & Anual                                             & 0,5 \textmugreek g/m3                                                                 \\ \hline
		\ce{C6H6}                                  & Anual                                             & 5 \textmugreek g/m3                                                                   \\ \hline
		\ce{CO}                                    & Horario                                           & 330 \textmugreek g/m3                                                                 \\ \hline
		\ce{PM2,5}                                 & Anual                                             & 25 \textmugreek g/m3                                                                  \\ \hline
	\end{tabular}
		\caption{Tabla valores límite \cite{informe_2018}}
\end{table}

Si nos centramos en el temas únicamente referidos a la salud existen unos \textbf{valores objetivo} objetivo que se recogen en la siguiente tabla: 


\begin{table}[H]
	\begin{center}
	\begin{tabular}{|l|l|l|}
		\hline
		\textbf{Contaminante a analizar} & \textbf{Descripción} & \textbf{Valor Objetivo} \\ \hline
		\ce{PM2,5    }                        & Anual                & 25 \textmugreek g/m3                \\ \hline
		\ce{As           }                    & Anual                & 6 ng/m3                 \\ \hline
		\ce{Cd        }                       & Anual                & 5 ng/m3                 \\ \hline
	\ce{	Ni        }                       & Anual                & 20 ng/m3                \\ \hline
		\ce{B(a)P}                          & Anual                & 1 ng/m3                 \\ \hline
	\end{tabular}
\caption{Tabla valores objetivo \cite{informe_2018}}
	\end{center}
\end{table}

 \subsection{La calidad del aire en Madrid}
 
 Ahora hablamos tan solo de la comunidad de Madrid, que posee las siguiente características.
 
 % Please add the following required packages to your document preamble:
 % \usepackage{multirow}
 \begin{table}[H]
 		\begin{center}
 	\begin{tabular}{|c|c|c|c|}
 		\hline
 		\multicolumn{2}{|l|}{\textbf{Características}} & \multicolumn{1}{l|}{\textbf{Madrid}} & \multicolumn{1}{l|}{\textbf{España}} \\ \hline
 		\multirow{2}{*}{Población}       & (Habs.)     & 3.354.745                            & 46.722.980                           \\ \cline{2-4} 
 		& (\%)        & 7,18 \%                              & 100 \%                               \\ \hline
 		\multirow{2}{*}{Superficie}      & (km2)       & 7.424                                & 505.990                              \\ \cline{2-4} 
 		& (\%)        & 1,47 \%                              & 100 \%                               \\ \hline
 	\end{tabular}
 \caption{Tabla información Madrid. Fuente: INE \cite{informe_2018}}
 	\end{center}
 \end{table}
 
  Como conclusiones de los informes del pasado año se destaca principalmente lo siguiente:
 
 \begin{itemize}
 	\item En el año 2018 \textbf{se ha superado el valor límite anual de \ce{NO2}}, así como el valor objetivo de \ce{O3} tanto para la protección de la salud como de la vegetación.	Las causas de la superación del \ce{NO2} se atribuyen principalmente al tráfico de vehículos d	combustión ya que se trata de ubicaciones muy influenciadas por vías principales de tráfico.
 	\item En el ámbito de esta red no se supera el valor límite horario de \ce{NO2}, pero sí se produce, sin embargo, una \textbf{superación del valor límite anual de \ce{NO2}}, concretamente en la zona ES1308 “Corredor del Henares”, como consecuencia de los niveles alcanzados en la estación ES1869A “Coslada”, de tipo urbana de tráfico (41 \textmugreek g/m3 de media anual).
 \end{itemize}


    \imagen{mapaEspanaNO2}{\footnotesize{Mapa España superación límite legal durante 2018 con Madrid en rojo.} Fuente: Ecologistas en acción \cite{informe_ecologistas}}

Además de los valores emitidos en los informes anuales por el gobierno y otras organizaciones existen también numerosos papers \cite{nunez-alonso_statistical_2019} \cite{gomez-losada_data_2019} que estudian la creciente contaminación en la ciudad de Madrid.
 
 Madrid posee un un sistema de vigilancia que dispone de una red formada por 24 estaciones que pueden clasificarse en tres categorías en cuanto al tipo de ambiente en el que se ubican: 9 estaciones de tráfico (situadas próximas al viario), 12 estaciones de fondo urbano (más alejadas del tráfico, generalmente en parques) y 3 estaciones suburbanas (situadas fuera del núcleo urbano consolidado).
 
     \imagen{cartdoDBestaciones}{\footnotesize{Localización estaciones Madrid.} Fuente: Elaboración propia mediante \href{https://carto.com}{cartodb}
 	y los datos de madrid \cite{portal_datosabiertos_madrid}}

 Así, para el caso de los contaminantes que se analizan, tenemos que el \ce{NO2} se mide en las 24 estaciones, las partículas \ce{PM10} en 12 de ellas, las partículas \ce{PM2.5} en 6 y el \ce{O3} se registra en 14 estaciones. Por otro lado, el Ayuntamiento ha establecido una zonificación de la ciudad de Madrid orientada a la gestión de situaciones de altos niveles de contaminación, como los picos de contaminación por \ce{NO2}, que ponen en marcha la aplicación del protocolo aprobado por el Ayuntamiento de Madrid para hacer frente a dichas situaciones.
 
 Como medidas del protocolo de actuación cabe destacar que el dióxido de nitrógeno \ce{NO2} \textbf{nunca deberá sobrepasar los 200 \textmugreek g/m3}, de hacerlo se entraría en estado de ‘preaviso’.
 
 La legislación europea, a la cual se ajusta la ciudad de Madrid, establece también un valor límite horario de \ce{NO2}, para proteger a la población de exposiciones a altos niveles de este contaminante, aunque sea por cortos periodos de tiempo (denominados "pico de contaminación"). El valor límite horario para el \ce{NO2} está establecido en 200 \textmugreek g/m3, límite que no debería rebasarse más de 18 horas al año en ninguna estación de la ciudad.
 
    Madrid, por ser una de las ciudades más susceptibles de sufrir contaminación, tiene varios protocolos de actuación ante distintos niveles de contaminación del aire. Todos ellos se atienden al ICA (Índice de Calidad del Aire. Ver figura ~\ref{fig:tablaICA}). Mediante este índice se indica a si el aire es apto o no, o qué grado de contaminación tiene. 
  
 Según el mencionado índice, si el valor está entre \textit{0 y 50}, las \textbf{condiciones del aire son buenas}. Si se encuentra entre \textit{51 y 100}, son \textbf{regulares}. A partir de \textit{101}, y hasta el valor \textit{150}, el nivel de contaminación es \textbf{dañino} para la salud de algunos grupos (niños y ancianos, entre otros). Desde los valores \textit{151 hasta 200}, \textbf{el aire es contaminante para cualquiera}. \textit{A partir de 201}, los niveles son \textbf{muy dañinos e incluso peligrosos}.

En cuanto a las  situaciones de aviso existen varias: 
\begin{itemize}
	\item \underline{Preaviso}: cuando dos estaciones cualesquiera detectan un nivel superior a 180 microgramos/m3 durante dos horas consecutivas.
	\item \underline{Aviso}: cuando dos estaciones cualesquiera detectan un nivel superior a 200 microgramos/m3 durante dos horas consecutivas
	\item \underline{Alerta}: cuando tres estaciones cualesquiera detecten un nivel superior a 400 microgramos/m3 durante tres horas consecutivas.
\end{itemize}

Durante años anteriores se han desarrollado diferentes protocolos que van en función del tipo de situaciones de aviso y de la legislación que el gobierno vigente acometa o apruebe en cada legislatura. Una vez se haya superado o se prevea superar alguno de los niveles citados en las situaciones de aviso, y si la previsión meteorológica es desfavorable, se considerará iniciado un episodio de contaminación. En ese contexto se dan varios escerarios posibles con diferentes actuaciones.


\begin{itemize}
	\item \underline{ESCENARIO 1}
		\begin{itemize}
			\item[$\ast$] \textbf{1 día con superación del nivel de preaviso. }
			\\ \\
			Actuaciones:
		\end{itemize}
		\begin{itemize}
				\item Medidas informativas y de recomendación.
				\item Medidas de promoción del transporte público.
				\item Reducción de la velocidad a 70 km/h en la M-30 y accesos.
		\end{itemize}
	\item \underline{ESCENARIO 2}
		\begin{itemize}
			\item[$\ast$] \textbf{2 días consecutivos con superación del nivel de preaviso o 1 día con superación del nivel de aviso. }
			\\ \\
			Actuaciones:
		\end{itemize}
		\begin{itemize}
			\item Todas las medidas del escenario 1.
			\item Prohibición de la circulación en el interior de la M-30 y por la M-30 a los vehículos a motor, incluidos ciclomotores, que no tengan la clasificación ambiental de “CERO EMISIONES”, “ECO”, “C” o “B” en el Registro de Vehículos de la Dirección General de Tráfico. 
			\item Prohibición del estacionamiento en las plazas y horario del Servicio de Estacionamiento Regulado (SER) a los vehículos a motor que no tengan la clasificación ambiental de “CERO EMISIONES” o “ECO” en el Registro de Vehículos de la Dirección General de Tráfico. 
		\end{itemize}
	\item \underline{ESCENARIO 3}
		\begin{itemize}
			\item[$\ast$] \textbf{3 días consecutivos con superación del nivel de preaviso o 2 días consecutivos con superación del nivel de aviso. }
			\\ \\
			Actuaciones:
		\end{itemize}
		\begin{itemize}
			\item Todas las medidas del escenario 1.
			\item Prohibición del estacionamiento en las plazas y horario del Servicio de Estacionamiento Regulado (SER) a los vehículos a motor que no tengan la clasificación ambiental de “CERO EMISIONES” o “ECO” en el Registro de Vehículos de la Dirección General de Tráfico. 		
			\item Se recomienda la no circulación de taxis libres, excepto Eurotaxis y vehículos que tengan la clasificación ambiental de “CERO EMISIONES” o “ECO” en el Registro de Vehículos de la Dirección General de Tráfico en todo el término municipal. Estos vehículos podrán estacionar en las plazas del SER, además de en sus paradas habituales a la espera de viajeros, en los términos que se establezcan en la Ordenanza de Movilidad Sostenible.
		\end{itemize}
	\item \underline{ESCENARIO 4}
		\begin{itemize}
			\item[$\ast$] \textbf{4 días consecutivos con superación del nivel de aviso. }
			\\ \\
			Actuaciones:
		\end{itemize}
		\begin{itemize}
			\item Todas las medidas del escenario 1.
			\item Prohibición del estacionamiento en las plazas y horario del Servicio de Estacionamiento Regulado (SER) a los vehículos a motor que no tengan la clasificación ambiental de “CERO EMISIONES” o “ECO” en el Registro de Vehículos de la Dirección General de Tráfico. 	
			\item Se recomienda la no circulación de taxis libres, excepto Eurotaxis y vehículos que tengan la clasificación ambiental de “CERO EMISIONES” o “ECO” en el Registro de Vehículos de la Dirección General de Tráfico en todo el término municipal. Estos vehículos podrán estacionar en las plazas del SER, además de en sus paradas habituales a la espera de viajeros, en los términos que se establezcan en la Ordenanza de Movilidad Sostenible.		
			\item Prohibición de la circulación en el interior de la M-30 y por la M-30 a los vehículos a motor, incluidos ciclomotores, que no tengan la clasificación ambiental de “CERO EMISIONES”, “ECO” o “C” en el Registro de Vehículos de la Dirección General de Tráfico. 
		\end{itemize}
\end{itemize}

    \imagen{protocolo}{\footnotesize{Infografía escenarios.} Fuente:  \href{https://www.madrid.es/}{Madrid/contaminación}}

Como vemos son medidas acumulativas que dependen de la gravedad de la situación. Este tipo de medidas no solo permite hacer descender el nivel de contaminación durante los días que están activas, sin que también incentiva la compra de vehículos híbridos o vehículos totalmente eléctricos. 

Debido a los diversos toques de atención de la Unión Europea Madrid cuenta con un protocolo de actuación claro y eficaz del que carecen otras ciudades españolas.

 \subsection{Low Emission Zone (LEZ): Madrid  Central}
También denominadas ZUAP (Zonas Urbanas de Atmósfera Protegida \cite{wayback_2012}) o ZBE (Zona de bajas emisiones), son aquellos espacio dentro de una ciudad que tienen vetada la entrada a los vehículos mas contaminantes al espacio delimitado, en algunos casos la ZBE, solo incluye aquellos vehículos denominados “pesados o industriales” y en otros, también a todos los que entren dentro de la zona. Algunos de los países que implementan, desde hace varios años, este tipo de medidas son Noruega, Francia, Holanda o Gran Bretaña. 

Madrid Central es una zona de emisiones que  entró en vigor el 30 de noviembre del 2018, esta medida está incluida dentro del Plan A de Calidad del Aire y Cambio Climático, la medida tiene como pretensión principal, la reducción de los gases nocivos, bajando así un 40\% las emisiones en el distrito central de la ciudad y un 20\% los desplazamientos dentro de la zona.

La zona afectada está conformada por la mayoría de calles del centro de la Ciudad, entre las mas relevantes incluidas se pueden destacar las siguientes: Alberto Aguilera, Glorieta de Ruíz Jiménez, Carranza, Glorieta de Bilbao, Sagasta, Plaza de Alonso Martínez, Génova, Plaza de Colón, Paseo de Recoletos, Plaza de Cibeles, Paseo del Prado, Plaza de Cánovas del Castillo, Paseo del Prado, Plaza del Emperador Carlos V, Ronda de Atocha, Ronda de Valencia, Glorieta de Embajadores, Ronda de Toledo, Glorieta de la Puerta de Toledo, Ronda de Segovia, Cuesta de la Vega, Calle Mayor, Calle Bailén, Plaza de España (lateral continuación de la Cuesta de San Vicente), Calle Princesa y Calle Serrano Jover.

    \imagen{madridcentral}{\footnotesize{Perímetro Madrid Central.} Fuente:  \href{https://www-s.munimadrid.es}{Ayto Madrid}}

 \hl{Si tengo tiempo puedo hablar mas sobre Madrid Central en la parte de experimentos}

 \subsection{La calidad del aire y la salud}

El 95\% de la población mundial vive en áreas que no cumplen las pautas de un aire sano, según el informe State of Global Air \cite{state_global_air_2018}, del Health Effects Institute. Las ciudades, donde viven más de la mitad de los casi 7.500 millones de habitantes del planeta, son el caldo de cultivo de la contaminación, un importante factor de mortalidad, al que solo superan la hipertensión, la dieta no saludable y el tabaco.

El aire contaminado puede generar problemas de pulmón o del corazón, además de problemas disfuncionales cardio respiratorios e incluso la muerte, pero esto sería únicamente ante una exposición muy prolongada, o antes varias exposiciones periódicas prolongadas. De hecho, la OMS (Organización Mundial de la Salud) ha advertido que representa un importante riesgo medioambiental para la salud y que provoca cada año unas tres millones de defunciones prematuras, de las cuales medio millón corresponderían a Europa \cite{2018_oms_calidad}.

Según la Organización Mundial de la Salud (OMS), la contaminación del aire es actualmente uno de los mayores riesgos sanitarios mundiales, comparable a los riesgos relacionados con el tabaco  \cite{2018_oms_calidad}. Algunos agravios directos que puede producir serían los siguientes.

\begin{itemize}
\item Limitar las vías respiratorias.

\item Agravar o incluso generar enfermedades respiratorias.

\item Dañar partes profundas de los pulmones, aun después de que desaparecen ciertos síntomas como tos o dolor de garganta.

\item Riesgo mayor de padecer enfermedades cardiovasculares.

\item Algunos de los efectos de la exposición a niveles altos de monóxido de carbono pueden disminuir los reflejos y causar confusión y somnolencia.	
	
\end{itemize}



\section{Análisis legislativo (En España)}\label{teorico-legislativo}

Por supuesto toda esta maraña de terminología científica está regida por diversas leyes, tanto a nivel nacional como también Europeo. La presencia en la atmósfera de sustancias contaminantes, que pueden ser gases, partículas y/o aerosoles es la que determina en última instancia la calidad del aire. En España, la protección de la atmósfera y de la calidad del aire pasa por la prevención, vigilancia y reducción de los efectos nocivos de dichas sustancias contaminantes sobre la salud y el medio ambiente en su conjunto, en todo el territorio nacional. Para ello, la normativa vigente en materia de calidad del aire establece unos objetivos de calidad del aire, o niveles (concentraciones) de contaminantes en la atmósfera que no deben sobrepasarse. 

España comunica anualmente información sobre calidad del aire a la Comisión Europea en cumplimiento de diferentes directivas . En la siguiente lista se realiza un recopilatorio de las más importantes a nivel europeo y nacional. En el informe anual de 2018 \cite{informe_2018}, apartado 2.1 se añade información detallada sobre el marco legislativo y cada una de las leyes que incumber a nuestro país en materia de gases nocivos.

\begin{itemize}
\item  \underline{Marco legislativo Europeo}
	\begin{itemize}
	\item \textit{\textbf{Directiva 2008/50/CE} del Parlamento Europeo y del Consejo, de 21 de mayo de 2008, relativa a la calidad del aire ambiente y a una atmósfera más limpia en Europa.}
	\item \textit{\textbf{Directiva 2004/107/CE} del Parlamento Europeo y del Consejo, de 15 de diciembre de 2004, relativa al arsénico, el cadmio, el mercurio, el níquel y los hidrocarburos aromáticos policíclicos}
	en el aire ambiente.
	\item \textit{\textbf{Directiva 2015/1480/UE}, de la Comisión, de 28 de agosto de 2015, por la que se modifican varios anexos de las Directivas 2004/107/CE y 2008/50/CE del Parlamento Europeo y del Consejo en los que se establecen las normas relativas a los métodos de referencia, la validación de datos y la ubicación de los puntos de muestreo para la evaluación de la calidad del aire ambiente.}
	\item \textit{\textbf{Decisión de ejecución de la Comisión 2011/850/UE, de 12 de diciembre de 2011}, por la que se establecen disposiciones para las Directivas 2004/107/CE y 2008/50/CE del Parlamento Europeo y del Consejo en relación con el intercambio recíproco de información y la notificación sobre la calidad del aire ambiente.}
	\end{itemize}
\item  \underline{Marco legislativo Nacional}
	\begin{itemize}
	\item \textit{\textbf{Ley 34/2007}, de 15 de noviembre, de calidad del aire y protección de la atmósfera.}
	\item \textit{\textbf{Real Decreto 102/2011}, de 28 de enero, relativo a la mejora de la calidad del aire.}
	\item \textit{\textbf{Orden TEC/351/2019}, de 18 de marzo, por la que se aprueba el Índice Nacional de Calidad del Aire.}
	\end{itemize}
\end{itemize}

A modo de resumen, en la actualidad, los textos legales más relevantes para la calidad del aire en España son: la \textit{Directiva europea 2008/50}; la \textit{Ley 34/2007}, de Calidad del Aire y Protección de la Atmósfera; y el \textit{R.D. 102/2011 }relativo a la mejora de la calidad del aire.

Si nos atenemos únicamente a Madrid, que es la ciudad objeto de estudio en el presente trabajo, a parte de seguir todas estás leyes nacionales y europeas poseen un protocolo de contaminación propio para actuar en caso de sobrepasar los niveles establecidos en dichas leyes. No es frecuente que cambie pero si es interesante preguntar, debido a los cambios de gobierno, cuál es el que se encuentra en vigor. Consulté con la administración actual para recibir esa información y me respondieron amablemente. El protocolo de actuación se puede consultar de manera online \cite{ley_protocolo_madrid} por cualquier ciudadano.

    \imagen{protocoloRespuestaCorreo}{\footnotesize{Respuesta ayundamiento de Madrid.} Fuente: Consulta realizada \href{https://www-s.munimadrid.es}{Ayto Madrid}}

Para ampliar la información sobre normativa se puede consultar el apartado \textit{Acerca de Datos Abiertos / Normativa} en la \href{https://datos.madrid.es/portal/site/egob/menuitem.400a817358ce98c34e937436a8a409a0/?vgnextoid=830512b9ace9f310VgnVCM100000171f5a0aRCRD&vgnextchannel=830512b9ace9f310VgnVCM100000171f5a0aRCRD&vgnextfmt=default}{página web} de datos abiertos de la comunidad de Madrid que es el mismo que aplica a toda España.

Para terminar, y para conocer por qué me he decantado finalmente por unos contaminantes y no otros a la hora de realizar el estudio, en el \textit{Real Decreto 102/2011} relativo a la mejora de la calidad del aire, establece umbrales de alerta para tres contaminantes,  dióxido de nitrógeno (\ce{NO2}), dióxido de azufre (\ce{SO2}) y ozono (\ce{O3}). Define además el umbral de alerta como el nivel a partir del cual una exposición de breve duración supone un riesgo para la salud humana, que afecta al conjunto de la población y que requiere la adopción de medidas inmediatas. El valor del umbral de alerta para el dióxido de nitrógeno está establecido en \textbf{400 \textmugreek g /m3} durante tres horas consecutivas en lugares representativos de la calidad del aire, en un área de al menos 100 km2 o en una zona o aglomeración entera, si ésta última superficie es menor..

\section{Análisis parte técnica}\label{teorico-tecnico}

En esta parte vamos a realizar una inmersión a la parte teórica de la parte más técnica del proyecto. Servirá para conocer un poco más a fondo la parte tecnológica con la que hemos realizado todo el trabajo.

\subsection{Fuente de  datos}

Los datos provienen de \href{https://datos.madrid.es}{datos Madrid} que es el portal de datos abiertos de la comunidad de Madrid. Tal y como ellos mismos se describen en la web: \textit{El portal de datos abiertos de la comunidad de Madrid contiene  datos en bruto (bases de datos) que pueden ser utilizados, reutilizados y redistribuidos libremente por cualquier persona, y que ponemos las administraciones a disposición de la ciudadanía para realizar nuevas aplicaciones, estudios, análisis, investigaciones, etc.}

La verdad que la página está realmente bien, hay \textit{datasets} que contienen informaciones bastante relevantes, se actualizan de manera frecuente y las opciones de descarga son múltiples lo que permite versatilidad a la hora de realizar proyectos. Aunque se realizará un análisis exhaustivo más adelante de la tipología de las fuentes de origen, es interesante el saber por qué se ha escogido\textit{ .csv} como archivo de datos de las fuentes. 

   \imagen{tipodatos}{\footnotesize{Tipos de datos ofrecidos por Ayto. Madrid.} Fuente:  \href{https://datos.madrid.es}{Datos Madrid}}

Un csv \textit{(comma-separated values}) es un archivo de texto que almacena los datos en forma de columnas, separadas generalmente por coma y donde las filas se distinguen por saltos de línea. Este tipo de archivos podría decirse que es el lenguaje internacional de trabajo con datos, aunque por supuesto no es el único \cite{formatos}. De hecho tiene un estándar \cite{estandar} por el que se rige, si bien es cierto que cada persona lo interpreta un poco como conviene para cada proyecto.
En el caso concreto del estudio yo he descargado los datos en formato \textit{.csv}, principalmente los datos de la calidad del aire, los datos de las estaciones y los datos de temperatura. 

\subsection{Modelo analítico}
Alrededor de un 80\% del análisis de un proyecto o negocio se invierte en los procesos iniciales, es decir, en la elección de un buen modelo que permitar extraer, transformar y cargar los datos de una manera óptima. 

Para ello lo más importante que tenemos que tener en cuenta antes de abordar un proceso que involucre grandes masas de datos a analizar es el objetivo que se persigue. Responder a preguntas como\textit{ ¿qué queremos saber?}, \textit{¿qué datos son necesarios para saberlo?}, \textit{¿cómo se relacionan esos datos entre sí?} son algunas de las preguntas básicas que nos debemos hacer. La razón de ser de cada modelo analítico es el objetivo del proyecto para el que se trabaja. 

Existen modelos de varios tipos: \textbf{predictivos}, basado en la premisa de si ocurre X pasará Y, \textbf{prescriptivos}, ayudan a facilitar la toma de decisiones, y \textbf{descriptivos}, utilizados entre otras cosas para establecer relaciones entre los datos. Al final lo que se intenta con todos ellos es construir sistemas que encuentren patrones útiles en los datos que consigan responder a las hipótesis iniciales planteadas de cada proyecto.

A continuación se resume el procedimiento de manera genérica. 
\begin{enumerate}
	\item El proceso de análisis comienza con la recopilación de datos, en la cual se identifican la información que necesitan para una aplicación de análisis en particular. Es posible que sea necesario combinar los datos de diferentes sistemas de origen a través de rutinas de integración de datos, transformarlos en un formato común y cargarlos en un sistema de análisis. En otros casos, el proceso de recopilación puede consistir en extraer un subconjunto relevante de una secuencia de datos brutos y moverlo a una partición separada en el sistema para que pueda analizarse sin afectar el conjunto de datos completo.
	\item Una vez que los datos que se necesitan están guardados, el siguiente paso es encontrar y corregir los problemas de calidad de los datos que podrían afectar la precisión de las aplicaciones de análisis. Eso incluye la ejecución de perfiles de datos y tareas de limpieza de datos para garantizar que la información en un conjunto sea coherente y que se eliminen los errores y las entradas duplicadas.
	\item Luego, se realizan trabajos adicionales de preparación para manipular y organizar los datos para el uso analítico planificado, y se aplican políticas de gobierno de datos para garantizar que los datos se ajusten a los estándares corporativos y se usen correctamente.
\end{enumerate}

En otra palabras se podría resumir con lo que se conoce como \textit{proceso ETL}: \textbf{Extracción} (obtención de datos de las fuentes de origen), \textbf{Transformación} (realización de los cálculos necesarios para obtener los datos que nos interesan) y \textbf{Carga} (en esta parte del proceso se vuelcan los datos procedentes de la fase de transformación al sistema de destino.)

Finalmente hay que comunicar los resultados. Una manera idónea son las visualizaciones de datos.

 \hl{Aqui puedo darle una vuelta e incluso hacerlo mas detallado}

\subsection{Visualización de datos}
Una vez tenemos claro el modelo con el que se van a obtener los datos la última parte comprende la visualización de los mismos.

Desde un punto de vista más comercial podríamos afirmar que la visualización de datos es un proceso que consume datos como entrada y los transforma en conocimiento del negocio.

Desde un punto de vista más técnico la visualización de datos es la presentación gráfica de información con dos propósitos principales. Por un lado, la interpretación y construcción de significado a partir de los datos (es decir, el análisis); y por otro lado, la comunicación (la comunicación de la información que se obtiene de esos datos). De una manera más simple se trata de representar magnitudes de forma visual con el objetivo último de presentar una información. Normalmente esta información se presenta a una audiencia o puede ser también parte de un estudio con el que obtener patrones o modelos dentro de un proyecto más grande. 

   \imagen{dataviz}{\footnotesize{\textit{What makes a good visualization?} by David McCandless.} Fuente:  \href{https://InformationIsBeautiful.net}{InformationIsBeautiful.net}}

Una  visualización es una herramienta muy potente para descubrir y comprender la lógica que se encuentra detrás de un conjunto de datos, así como para compartir esta interpretación con otras personas desde un punto de vista objetivo.

El potencial de la visualización de datos es muy muy alto. En muchas ocasiones solemos tener una maraña de datos que por sí solos no nos dicen nada, sin embargo, una vez procesados, transformados y limpios, unos datos pueden convertirse a través de una buena visualización en una información muy relevante.

Algunas organizaciones \cite{cesal_what_2019} tiene sus propias guías de estilo para la creación de visualizaciones. Es el caso de por ejemplo Google que tiene un apartado de \href{https://material.io/design/communication/data-visualization.html#types}{Material Design} dedicado en exclusiva al \textit{data visualization}.

\capitulo{4}{Técnicas y herramientas}

Esta parte de la memoria tiene como objetivo presentar las tecnologías y las herramientas de desarrollo que se han utilizado para llevar a cabo el proyecto. Se han estudiado diferentes alternativas de metodologías, herramientas, bibliotecas  y se pretende aquí realizar un resumen de los aspectos más destacados, incluyendo comparativas entre las distintas opciones y, en caso de ser necesario, una pequeña justificación de las elecciones realizadas. 

No se pretende que este apartado se convierta en un capítulo de un libro dedicado a cada una de las opciones de manera detallada, sino comentar de manera breve los aspectos más destacados de cada técnica o herramienta que ha sido empleada.

\section{Metodologías}\label{metodologias}
En este apartado se describen las metodologías utilizadas para el desarrollo del proyecto, concretamente a la parte diaria de desarrollo de producto y a la parte de como hemos enfocado la investigación.

\subsection{Estrategia de investigación}
En el presente proyecto se sigue una estrategia de investigación de análisis de datos cuantitativos. Y dentro de ésta se busca realizar una anális exploratorio y descriptivo que de pie a sacar una conclusiones de los datos escogidos como punto de partida

\begin{quote}
\textit{"La idea del análisis de datos es buscar patrones en los datos y sacar conclusiones. Existe una amplia gama de técnicas establecidas para analizar datos cuantitativos"} \cite{oates_researching_2006}. 
\end{quote}

Siguiendo esta estrategia se utilizan técnicas estadísticas exploratorias y descriptivas simples para encontrar patrones en los datos y comprobar si dichos patrones se encuentran realmente en los datos o si son sólo fruto del azar. Además, se utilizan tablas, mapas y gráficos para presentar los datos de una manera visual y sencilla. La recolección de datos se realiza mediante portales de datos abiertos vía descarga ordinaria. 

En cuanto a la realización del proyecto respecto de la parte técnica, se hace uso de \textit{Python} como lenguaje de programación de uso general y de \textit{R} como entorno y lenguaje de programación con enfoque estadístico. Además también he usado \textit{SQL} para realizar diversas pruebas. Para los gráficos me apoyo en diferentes herramientas \textit{open source} y par la visualización final en \textit{PowerBi}. También se realizará una pequeña web a modo de resumen para recopilar información relevante del proyecto. 

\subsection{Metodología de trabajo diaria: Scrum}\label{metodologias_scrum}
Scrum es un sistema de desarrollo de software que está dentro de las metodologías ágiles y es una de las más conocidas actualmente. A grandes rasgos podemos decir que se basa en la creación y asignación de tareas. Esta tarea entrará en un sistema iterativo que controlará los cambios de estado de la tarea, hasta que esta resulta completada y se descarta del panel de \textit{Scrum}. 

Se trata de una herramienta muy útil en espacios donde los grupos de trabajo tienen dificultades para hacer las acciones u operaciones que les lleven a objetivos en común. Dicho de otro modo, \textit{Scrum} sirve para que equipos multidisciplinares trabajen en entornos complejos, donde los requisitos son muy cambiantes, y los resultados se tienen que obtener en un plazo corto de tiempo. 

\imagen{ScrumSprint}{\footnotesize{Resumen metodología \emph{Scrum}. Fuente: \url{manifesto.co.uk}}}

Podemos destacar tres pilares básicos en \textit{Scrum}:

\begin{itemize}
	\item Transparencia: Todos los aspectos del proceso son visibles para todos aquellos que administran el proyecto.
	\item Inspección: Se debe controlar con la frecuencia suficiente los diversos aspectos del proceso para que puedan detectarse variaciones inaceptables en el mismo. Un ejemplo serían las \textit{dailys}, una reunión diaria corta, de no más de quince minutos, donde cada miembro del equipo expone el trabajo que está realizando.
	\item Revisión: El producto debe estar dentro de los límites aceptables por el cliente. Una buena metodología necesita de un mecanismo de mejora continua, es decir, de control, para adaptarse y mejorar en caso de ser necesario.
	
\end{itemize}

En resumen, Scrum propone seguir un proceso de desarrollo iterativo e incremental a través de una serie de iteraciones denominadas \textit{sprints} y de revisiones. 
Hablando en términos personales, obtuve el \textit{título certificado de Scrum Master} el pasado mes de abril y es una metodología que aplico a diario durante el trabajo y que particularmente me apasiona. Si bien es cierto, siguiendo la teoría, que es necesario tener en cuenta que esta metodología esta pensada para trabajar en equipo por lo que en este trabajo se ha intentado emular de la mejor forma posible dadas las circunstancias: seguimiento de un \emph{board}, reuniones puntuales con los tutores, \emph{sprints} quincenales\ldots

\imagen{BoardExample1}{\footnotesize{\emph{Board} de tareas empleado durante el proyecto. Fuente:  Elaboración propia.}}

Se detallará más la metodología que se ha seguido en este proyecto en el \textbf{Apéndice A}.
\section{Herramientas}\label{herramientas}
En este apartado se describen las herramientas utilizadas para el desarrollo del sistema, concretamente en esta parte hago referencia a la parte software.

\subsection{Pycharm}\label{herramientas_pycharm}
\href{https://www.jetbrains.com/pycharm/}{Pycharm} es el editor es un IDE (Entorno de desarrollo integrado) desarrollado por la compañía Jetbrains, está basado en\textit{ IntelliJ IDEA}. Pycharm tiene cientos de funciones que lo puede ver como una herramienta muy pesada, pero que valen la pena ya que ayuda con el desarrollo del día a día. Es el entorno profesional por excelencia para trabajar con \textit{Python}. Es posible configurar también\textit{ jupyter notebooks} para trabajar de maner interna con este editor.

\begin{itemize}
	\item \subsubsection{Alternativas estudiadas}
	\begin{itemize}
		\item \href{www.sublimetext.com}{Sublime Text}: se pensó en esta alternativa debido a que se trata de un IDE bastante liviano en comparación con Pycharm. Finalmente se descartó por aspectos personales.
	\end{itemize}
\end{itemize}


\subsection{Jupyter Notebooks}\label{herramientas_jupyter}
\href{https://www.jetbrains.com/pycharm/}{Jupyter Notebook} es un entorno de trabajo interactivo que permite desarrollar código en \textit{Python} de manera dinámica, a la vez que integrar en un mismo documento tanto bloques de código como texto, gráficas o imágenes. Es un \textit{SaaS} utilizado ampliamente en análisis numérico, estadística y \textit{machine learning}, entre otros campos de la informática y las matemáticas. 

\subsection{Conda}\label{herramientas_conda}
\href{https://docs.continuum.io/anaconda/navigator/getting-started/ }{Conda} (Anaconda Navigator) es una interfaz gráfica de usuario para el gestor de paquetes y entornos de conda. A través del cuál se ha instalado \textit{ Jupyter Notebooks}.

\subsection{R studio}\label{herramientas_r_studio}
\href{https://www.rstudio.com}{R studio} es un entorno de desarrollo integrado (IDE) para el lenguaje de programación R, dedicado a la computación estadística y gráficos. Incluye una consola, editor de sintaxis que apoya la ejecución de código, así como herramientas para el trazado, la depuración y la gestión del espacio de trabajo. Aunque existen otros, sin duda es el IDE preferido para trabajar con R.

\subsection{Visual Studio code}\label{herramientas_atom}
\href{https://code.visualstudio.com}{Visual Studio code} es un editor de código fuente desarrollado por Microsoft para Windows, Linux y macOS. Incluye soporte para depuración, control de Git embebido, resaltado de sintaxis, finalización inteligente de código, fragmentos y refactorización de código. También es personalizable, para que los usuarios puedan cambiar el tema del editor, los atajos de teclado y las preferencias. Es gratuito y de código abierto, aunque la descarga oficial se realiza bajo licencia propia. 

Lo mejor de este editor es la gran cantidad de \textit{plugins} que puedes instalar facilitando programar en casi cualquier lenguaje.

\imagen{vscode}{\footnotesize{Pantalla principal \textit{VSCODE}. Fuente: \url{https://code.visualstudio.com}}}

\subsection{Firefox for developers}\label{herramientas_firefox}
Para el navegador he elegido Firefox en su versión para desarrolladores: \href{https://www.mozilla.org/es-ES/firefox/developer}{Firefox Developer Edition}, el cuál considero que posee todas las herramientas esenciales para testear y probar una aplicación web.


\section{Herramientas de visualización de datos}\label{herramientasdataviz}
En este apartado se describen las herramientas utilizadas para el desarrollo la parte de visualización de datos.

\subsection{Power BI}
\href{https://powerbi.microsoft.com}{Power BI} es una solución de análisis empresarial que permite visualizar los datos y compartir información con toda la organización, o insertarla en su aplicación o sitio web. En este caso se ha utlizado a través de una máquina virtual windows proporcionada por la universidad.

\begin{itemize}
	\item \subsubsection{Alternativas estudiadas}
	\begin{itemize}
		\item \href{https://www.qlik.com/us}{Qlik}, fue visto y utilizado durante el máster y es sin duda una gran alternativa aunque me encuentro más cómodo trabajando con PowerBI.
		\item \href{https://www.tableau.com}{Tableau}, esta fue nuestra  segunda opción. Tableu es una herramienta de visualización muy potente que permite realizar visualizaciones dinámica muy completas. Finalmente fue descartado por problemas de hardware. 
	\end{itemize}
\end{itemize}

\subsection{CartoDB}
\href{https://carto.com}{CartoDB} es una de las empresas líderes en cartografía mundiales. Es un servicio de pago aunque actualemente se ha accedido a él mediante una licencia universitaria que posee github en la que se tiene acceso a diversas herramientas entre las que forma parte CartoDB. La verdad es una suerte poder usarla ya que es un software muy sencillo de usar. Ha sido utilizado para la distribución geográfica de las estaciones.

\begin{itemize}
	\item \subsubsection{Alternativas estudiadas}
	\begin{itemize}
		\item \href{https://www.mapbox.com/}{Mapbox}  es una herramienta open source que permite geolocalizar posiciones entre otras opciones.
	\end{itemize}
\end{itemize}

\subsection{Vega}
\href{http://vega.github.io/}{Vega} es una herramienta \textit{open source} que permite, a través de la ingesta de diferentes formatos, crear gráfico interactivos. 

\begin{itemize}
		\item \href{http://vega.github.io/}{Voyager2} es una herramienta interna de vega que, dado un dataset, permite realizar un análisis exploratorio de las variablas que lo forman a través de diversos gráficos.
\end{itemize}

\subsection{DataWrapper}
\href{https://www.datawrapper.de}{DataWrapper} es una herramienta que dado un dataset (generalmente ya preprocesado ya que no se trata de una herramiente de limpieza sino que se enfocan más en la parte de visualizado) permite generar gráficos, mapas o tablas con el objetivo principal de enriquecer contenido en historias escritas. Principalmente está pensado para el acompañamiento de parte gráfica a noticias en diferentes medios de comunicación.


\section{Tecnologías}\label{tecnologias}
 En este apartado se describen las tecnologías utilizadas, concretamente lenguajes de programación, para el desarrollo del proyecto. 
 
 \subsection{Python}\label{tecnologias_python}
 
 \href{python.org}{Python} es un lenguaje de programación interpretado cuya filosofía hace hincapié en una sintaxis que favorezca un código legible. Los programas escritos en\textit{ Python} no necesitan compilarse de antemano para poder ejecutarse, por lo que es fácil probar pequeños fragmentos de código y hacer que el código sea más fácil de mover entre las plataformas.
 
Este lenguaje tiene licencia de software libre y dispone de una gran comunidad de usuarios que se encarga de enriquecerlo mediante la creación de nuevas librerías, funciones, etc. 

En el ámbito de la minería de datos, existen dos motivos principales por los que Python es tan popular:  	
\begin{itemize}
	\item Es muy sencilla la integración con otras aplicaciones con otras aplicaciones
	\item Dispone de una gran cantidad de librerías que implementan funciones para cálculos matemáticos y estadísticos, entre otros.
	 \begin{itemize}
		\item Algunas de las librerías que se han empleado en este proyecto han sido:
		
		 \hl{ANADIR LIBRERIAS}
		\begin{itemize}
			\item \href{https://www.qlik.com/us}{ggplot} 
			\item \href{https://www.qlik.com/us}{ggplot} 
		\end{itemize}
\end{itemize}


\end{itemize}
  \subsection{R}\label{tecnologias_R}
  \href{r-project.org}{R} es un entorno y un lenguaje de programación enfocado en el análisis estadístico de los más utilizados en el campo de la minería de datos que pueden aplicarse a gran variedad de disciplinas. De nuevo vuelve ser de software libre y es uno de los lenguajes de programación más utilizados en investigación científica.
  
   \begin{itemize}
		\item Algunas de las librerías que se han empleado en este proyecto han sido:
		 \hl{ANADIR LIBRERIAS}
  	\begin{itemize}
  		\item \href{https://www.qlik.com/us}{ggplot} 
  	\end{itemize}
  \end{itemize}
  

 \subsection{Tinybird}
\href{https://tinybird.co//}{Tinybird.co} es un \textbf{aplicación en beta} que gracias a un cúmulo de circustancias hemos tenido la suerte de probar en este proyecto. Se trata de una aplicación que utiliza lenguaje \textit{SQL} para realizar consultas sobre datasets. Es decir, dado un \textit{dataset} se crean diversos \textit{pipelines}, esa es la denominación que le dan, con los que podemos analizar, transformar y limpiar los datos de acuerdo a nuestros objetivos. Parece una especie de \textit{notebook} de jupyter aunque no lo es dado que no posee tal cantidad de opciones. 

Una vez tenemos los datos preprocesados podemos establecer un \textit{endpoint} a través del cuál es posible acceder de manera externa mediante \textit{API REST}. Después es muy fácil consumirlo desde, por ejemplo, un apicación web llamando a ese endpoint y acediendo a los datos ya limpios y listos para visualizarlos. 

 \hl{Terminar de redactar. anadir imagenes}

\section{Documentación}\label{docs}
En este apartado se describen algunas de las herramientas utilizadas para la parte de la documentación del proyecto.

 \subsection{La\TeX}\label{docs_latex}
  La\TeX es un sistema de composición de textos, orientado a la creación de documentos escritos que presenten una alta calidad tipográfica. Por sus características y posibilidades, es usado de forma especialmente intensa en la generación de artículos y libros científicos. 
  
   \begin{itemize}
  	\item \subsubsection{Detalles}
  	\begin{itemize}
  		\item \href{https://www.texstudio.org}{\textit{TeXstudio}}: como editor de escritorio de La\TeX se ha utilizado \textit{TeXstudio}, una recomendación personal de Carlos, tutor académico, y que sin lugar a dudas a ha supuesto un gran salto de calidad a la hora de realizar el trabajo en cuanto a la parte de la memoria.
  		\item \href{http://www.tablesgenerator.com/latex_tables}{\textit{Tables\_generator}}: una de las cosas más engorrosas a la hora de usar La\TeX\  son la tablas. Por esa razón se ha empleado \textit{tables\_generator} para realizar esta parte de una manera más cómoda. 
  		\item \href{https://es.sharelatex.com/}{ShareLatex}  es un servicio online que nos permite crear y compartir documentos en La\TeX. La característica por la que he empleado este editor es por que  permite la edición conjunta de un proyecto desde el navegador y ofrece diferentes tipos de plantillas según el uso que le vayamos a dar. Por lo tanto para realizar tablas, pruebas de imágenes u otros ensayos antes de copiarlo al documento final se han realizado todas con ayuda de esta herramienta.
  	\end{itemize}
  \end{itemize}
  
 \subsection{Zotero}\label{docs_zotero}
 \href{https://www.zotero.org/}{Zotero} es un gestor de referencias bibliográficas, libre, abierto y gratuito desarrollado por el \textit{Center for History and New Media} de la Universidad George Mason que funciona también como servicio. El funcionamiento de zotero es sencillo y está basado en los principios: recopilar, organizar, citar, sincronizar y colaborar.
  
  \imagen{zotero}{\footnotesize{Pantalla principal Zotero. Fuente: elaboración propia}}
  
   \begin{itemize}
   	\item \subsubsection{Alternativas estudiadas}
   		\begin{itemize}
	 	\item \href{https://www.mendeley.com/}{Mendeley}  es una aplicación web y de escritorio, propietaria y gratuita. Permite gestionar y compartir referencias bibliográficas y documentos de investigación, encontrar nuevas referencias y documentos y colaborar en línea. Fue descartada debido a problemas en otros trabajos anteriores con la sincronización entre dispositivos.
	   	\item \href{https://www.zenhub.com/}{Zenhub} es otro gestor de referencias bibliográficas que descarté por experimentar problemas con la extensión del navegador.
		\end{itemize}
	\end{itemize}
  
  \section{Otras herramientas}\label{otrasherramientas}
En este apartado se describen otras herramientas que se han empleado también durante el proyecto.

 \subsection{Git}
Git \cite{git_2019} es un sistema de control de versiones. Este tipo de sistemas registran los cambios realizados sobre un archivo o conjunto de archivos a lo largo del tiempo, de modo que es posible recuperar versiones específicas del mismo archivo más adelante.

 \subsection{Github}
\href{https://github.com/}{Github} es una plataforma cloud de desarrollo colaborativo de software para alojar proyectos utilizando el sistema de control de versiones \textit{git} \cite{git_2019}. El código se almacena de forma pública, aunque también se puede hacer de forma privada, además desde principios de este mismo año las opciones de repositorios privados se permiten de manera gratuita.  GitHub aloja tu repositorio de código y te brinda herramientas muy útiles para el trabajo en equipo. 

Se ha utilizado esta herramienta a modo de control de versiones y también se ha usado el \textit{project board} interno como gestor de tareas y de \textit{sprints}. Los \textit{project boards} funcionan, generalmente por repositorio y se pueden utilizar para crear flujos de trabajo personalizados que se adapten a las necesidades de cada proyecto. Algunos ejemplos podrían ser el seguimiento de tareas, la priorización del trabajo de características específicas, hojas de ruta completas o incluso listas de comprobación de versiones. Las tareas, desarrollo de la metodología y comentarios acerca del uso del \textit{project board} se pueden ver de manera más detallada en la parte de desarrollo del proyecto.

Se ha escogido el \textit{board} interno de Github respecto a las alternativas ya que considero que ha mejorado bastante con el tiempo y resulta muy útil cuando también tienes el código alojado en la misma plataforma.

\begin{itemize}
	\item \subsubsection{Alternativas estudiadas}
	\begin{itemize}
		\item \href{https://www.zenhub.io/}{Zenhub} es una extensión de Chrome para github. Se utiliza para gestionar proyectos y funciona de manera nativa en la interfaz. Se basa en la metodología ágil y resulta verdaderamente útil a la hora de realizar y gestionar un proyecto.
		\item \href{https://www.trello.com/}{Trello} es un gestor de tareas que permite el trabajo de forma colaborativa mediante tableros compuestos de columnas  que representan distintos estados. Se basa en el método \textit{Kanban} para gestión de proyectos, con tarjetas que viajan por diferentes listas en función de su estado. 
	\end{itemize}
\end{itemize}

\subsection{Github pages}
Dentro de esas herramientas que se mencionaban al describir Github en el apartado anterior se encuentran \textbf{github pages}. Esta herramienta permite alojar sitios web estáticos gratuitamente dentro un repositorio de código. GitHub pages permite dos modalidades de publicación:
\begin{enumerate}
	\item La primera es “User site” (solo se podrá tener un sitio de este tipo por cuenta); en este caso el sitio web será publicado en username.github.io (siendo username el nombre de usuario de la cuenta).
	\item La segunda opción es “Project site” (proyectos ilimitados) el cual será publicado en \textit{username.github.io/repository} (siendo repository el nombre del repositorio).
\end{enumerate}

Es una manera sencilla, rápida, cómoda de realizar pruebas para un web. Para proyectos de carácter académico en la que mostrar unos resultados en una opción bastante recomendada.

\begin{itemize}
	\item \subsubsection{Alternativas estudiadas}
	\begin{itemize}
		\item Plataformas en la nube:  como por ejemplo como Amazon AWS u OpenShift son muy populares hoy en día y permiten prácticamente lo mismo que github pages.
		\item \href{https://www.heroku.com/}{Heroku}  es un servicio de almacenamiento en la nube que además tiene mecanismos y herramientas para que la puesta en producción de las aplicaciones web sea prácticamente automática.
	\end{itemize}
\end{itemize}
    
 \subsection{StackOverFlow}

\href{https://stackoverflow.com/}{Stack Over Flow} es una de las comunidades de desarrolladores más importantes del mundo en la que se responden cuestiones de diferentes lenguajes. Es una herramienta funalmental para programadores. Actualmente Python es el lenguajde programación con más consultas en la plataforma.

 \subsection{OpenRefine}

\href{https://github.com/OpenRefine/OpenRefine}{OpenRefine} es una potente herramienta \textit{open source} para trabajar con datos: limpiarlos, transformarlos de un formato a otro y ampliarlos con servicios web y/o datos externos. Se utilizó para categorizar y ver los datos en la parte de investigación.

  \imagen{openrefine}{\footnotesize{Detalle vista principal. Fuente: elaboración propia}}


\capitulo{5}{Aspectos relevantes del desarrollo del proyecto}

Este apartado pretende recoger de manera teórica y muy breve los aspectos más interesantes del desarrollo del proyecto así como su justificación: decisiones que se han tomado, avances evolutivos por tipo y problemas que surgieron durante toda la realización del proyecto.


\section{Elección del proyecto}\label{inicio-proyecto}
Cuando llegó la hora de buscar proyecto, entre las diferentes opciones existían proyectos de carácter más técnico, también orientados a temas mas puramente teóricos o de carácter más de investigación. Sin duda, ésta última fue la opción escogida por el reto que esto supone. No existe por lo tanto un ruta exacta a seguir en este tipo de trabajos, sino que se trata más de tener un objetivo concreto pero a sabiendas de que el camino nos va a conducir a sacar unos patrones de los datos que pueden ser sorprendente o incluso inesperados. 

\section{Formación}\label{formacion}

Toda la formación con la que se ha contado para realizar este trabajo se ha adquirido gracias al máster cuyo proyecto final representa este trabajo. En éste se han aprendido las bases suficientes como para realizar un trabajo de éstas características. Si bien es cierto que la motivación principal era conducir la carrera del alumno hacia el ámbito del \textit{data science} y no ha habido demasiada formación de \textit{R} durante el máster. Es por ello que para esta funcionalidad concreta se han realizado un par de cursos extra. Para concretar algo más los cursos que se han tomado han sido de la plataforma \href{https://openwebinars.net/}{Open Webinars}.

\section{Metodologías/estrategias aplicadas}\label{metodologia}

Ya se han nombrado las metodologías más usadas durante el proyecto: por un lado \textit{Scrum} como metodología de desarrollo iterativa y por otro la estrategia de investigación ha sido la de análisis de datos cuantitativos. 

La estrategia de investigación a permitido aliarse perfectamente a la metodología escogida, al desarrollo del código fuente y a la exploración de la parte científica y legislativa del proyecto.

En resumen, ambas decisiones han fomentado de manera positiva tanto la motivación como el progreso del proyecto.


\section{Desarrollo análisis datos}
El modelo empleado para el desarrollo del proyecto comprende la exploración, transformación y posterior visualización de datos, se detalla de manera precisa en el apartado 7. 

Hay que destacar aquí que cuando hablamos de \textit{Big Data} estamos hablando de cantidades realmente enormes pero en este caso, aún con todo los datos descargados de todos los años, no supone una cantidad excesivamente grande de datos. Estamos hablando de que podemos llegar a un poco más de millón de filas si juntamos todos los datos, por ejemplo para el fichero utilizado para la visualización final. Pero hablamos de \textit{31831} filas para el caso de los meses que llevamos de 2019 o \textit{54831} filas para todo 2018. 

Es por esta razón que se han elegido este tipo de lenguajes de programación, quizás más enfocados también a otras áreas pero perfectamente válidos para un estudio de este tipo.

 
\section{Desarrollo visualización}\label{grafico}
La parte de visualización comprende varias partes:

\begin{itemize}
	\item La parte final del \textit{dashboard}, realizada con \textit{PowerBI}.
	\item Todas las gráficas internas que se han realizado en los diferentes lenguajes de programación, presentes en esta memoria o en la web. La finalidad última de éstas gráficas es la de explorar los datos o poder hablar con cierto criterio sobre los datos .
	\item El uso de otras herramientas, ya nombradas, mediante las cuales también se han realizado gráficas. Una vez con los ficheros modificados es sencillo obtener gráficas.
	\item La  web a modo de visualización-resumen.
\end{itemize}

En general ha sido una parte costosa en términos de tiempo. Además se han intentado seguir guías de estilo \cite{visualisation_2016} adecuadas o tener en cuenta gamas de colores aptas realizando la visualización de la mejor manera posible. 

\section{Desarrollo web}\label{dweb}

Como complemento a la visualización se ha realizado también una pequeña web, alojada en \textit{github pages}, que pretende ser un pequeño resumen vertical de las partes que nos han parecido de mayor relevancia. Dado que algunas veces existen problemas de insertado a la hora de introducir páginas propias de \textit{PowerBi} se ha tomado esta decisión.

\section{Documentación}\label{docs}

La documentación se ha realizado de manera progresiva. Siempre que se podía se añadían partes de la misma  preguntando sobre mejoras a los tutores, que han ayudado convenientemente a la realización de la misma. Se han realizado dos envíos a los tutores antes del envío final para su corrección.

\section{Dificultades encontradas}\label{inicio-proyecto}

Una de las partes más importantes del proyecto, en cuestión de tiempo, ha sido la de familiarizarse con los conceptos científicos y legislativos que requería el mismo. Quizás es necesario algo de autocrítica en este punto dado que una parte importante del tiempo se ha ido aquí en desventaja de otras fases. 

Existen en los datos una gran falta de homogeneidad. A pesar de que la página del ayuntamiento de Madrid está bastante detallada y correctamente explicada los datos varían a lo largo de los años, no solo en formatos sino también en diferentes aspectos. 

Existen problemas como por ejemplo el renombrado de calles los cuáles se deben intuir por que no aparecen por ningún lado especificados. Éste último error se dió con el caso de \textit{estación Plaza Fernandez Ladreda} y \textit{estación Plaza Elíptica}, en algunas mediciones vienen determinadas como dos estaciones distintas. Sin embargo, en los datos para la interpretación de los ficheros de medición ambas tienen el código 56. Esto es debido a que sufrió un cambio de nombre. Si este tipo de cambios no se especifican al usuario que se encuentra analizando los datos resulta muy costoso en terminos de tiempo el averiguar donde está el fallo.

Respecto a los datos más actuales los de 2019, existen errores de magnitud como por ejemplo no existen los datos para \ce{PM10} en ninguna de las medidas. Esto es debido a un error que los técnicos deben subsanar, como ya se ha informado. Por otro lado hubiera sido interesante el comparar el cese de \textit{Madrid Central}, en junio de 2019, para comprobar si realmente ha afectado a los niveles totales. Lamentablemente actualmente solo se encuentra disponible un mes desde que la normativa finalizó. Esto es debido a que a fecha de presentación de este trabajo, septiembre de 2019, todavía no está disponible el mes de agosto. Además no existe una fecha fija en la que los técnicos cargan los datos en la web y se tardan en corregir los errores.

Algunos otros contratiempos encontrados han sido con \textit{Mapbox}. Desde aquí se ha intentado contactar con atención al cliente de este sotware, por el momento sin respuesta, ya que no esta funcionando como debiera. El problema radica en que los mapas básicos que ofrece \textit{PowerBi} no permite realizar mapas de calor, entre otros, que resultan muy útiles para observar la contaminación. \textit{Mapbox} es de pago pero permite una versión gratuita limitada por el número peticiones que se realizan cada hora a la \textit{API}. Actualmente la configuración, que se realiza mediante un \textit{token} de identificación privado, funciona a días. 
 

\section{Agradecimientos}\label{reconocimientos}
Este pequeño apartado es a modo de agradecimiento a las siguientes personas que han contribuido de manera directa o indirecta al proyecto.
\begin{itemize}
	\item \textit{Jorge Gomez}, \textit{Javi Santana} (fundadores de \textbf{tinybird}): además del acceso a la \textit{beta} de la aplicación se han prestado a responder de manera altruista todas las preguntas que se han realizado sobre la misma.
	\item \textit{José Amador Viejo} (Director general de control ambiental, \textbf{Ayto de Madrid}): ha respondido sin ningún problema a las dudas sobre legislación que se han realizado.
	\item \textit{Juan Barcena del Riego} (Portavoz de \textbf{Ecologistas en acción}): se han intercambiado varios correos, sobre todo para la parte legislativa y también de la parte ambiental.
	\item \textit{Javier Di Deco Sampedro} (Data Scientist en \textbf{Piperlab}): existe un \href{https://twitter.com/datoxnitro_bot}{bot de twitter} realizado por la mencionada empresa que publica en tiempo real los niveles de contaminación por dióxido de nitrógeno de la ciudad de Madrid. Se envió un mensaje a la empresa para aprender las metodologías empleadas en el desarrollo del \textit{bot} y Javier nos atendió de manera amable y eficaz.
	\item \textit{Brenda Valverde}, actual periodista de \textbf{Newtral} mantuvimos una interesante conversación de carácter periodístico sobre contaminación ambiental en general y Madrid central en particular de la que se obtuvieron diversas fuentes para recopilar información. Es posible seguir su cuenta de twitter \href{https://twitter.com/brbrebren}{aquí}.
	\item \textit{Carlos E.Vivaracho Pascual} (tutor académico, \textbf{Universidad de Valladolid}): siempre atento durante el desarrollo del proyecto y poniendo las cosas fáciles cuando había incertidumbre.
    \item \textit{Fernando Cuenca Cabezas} (tutor de empresa, \textbf{Minsait}): cómplice de algunas de las decisiones del proyecto siempre se ha prestado a la colaboración durante todo del desarrollo. Además ha resultado fundamental para abordar diversas cuestiones técnicas.
\end{itemize}


\capitulo{6}{Trabajos relacionados}

En esta sección se han recopilado algunos recursos  útiles que tratan sobre el mismo tema o temas de temáticas parecidas. Realmente hay muchos estudios sobre la calidad del aire. Quizás el hecho de tener los datos en modo \textit{open data} de a la gente más libertad para usar estos datos y realizar proyectos con ellos.

\section{Trabajos útiles}

\begin{itemize}
	\item \href{https://github.com/UlisesGascon/Charla-AireMadrid-los-datos-abiertos}{\textit{AireMadrid} 'La realidad de los datos abiertos'} (Ulises Gascón, 2017): Ulises nos habla de las dificultades que supone crear un proyecto \textit{Open Source} que utiliza los datos abiertos de calidad del aire que publica el Ayuntamiento de Madrid para informar y permitir la reutilización efectiva de esa información por parte de los usuarios y otros desarrolladores.
	\item \href{https://github.com/Toolson12/poisonMAD}{\textit{Respira Madrid}} (José M. Martínez, 2019): Implementación en Java para el tratamiento e interpretación de datos sobre la calidad del aire en Madrid.
	\item \href{https://github.com/chucheria/CalidadAire}{\textit{Estudio calidad aire Madrid}} (Beatriz@Chucheria, 2016): Análisis sobre la calidad del aire y predicción sobre los niveles de contaminación de la ciudad de Madrid.
	\item \href{https://www.rodrigodemiguel.es/AnalisisAireMadrid/}{\textit{Análisis de la contaminación y la calidad del aire en Madrid}} (Rodrigo de Miguel González, 2016): análisis detallado del \ce{NO2} del aire de la ciudad de Madrid, compuesto que utiliza el Ayuntamiento para poner en marcha los periodos de restricciones por alta contaminación, comprobando el resultado de dichos periodos a lo largo del año.
	\item \href{https://idus.us.es/xmlui/handle/11441/79499}{\textit{Big data aplicado al transporte y en las ciudades: adaptación a la ciudad de Sevilla}} (Alejandro Casado Reinaldos, Trabajo Fin de Grado, Universidad de Sevilla, 2018): estudio analítico de datos de la ciudad de Sevilla. 
	\item \href{http://oa.upm.es/47313/}{\textit{Análisis temporal multivariante de la contaminación atmosférica dentro del distrito de metropolitano de Quito}} (Juan Luis Manosalvas Paredes, Trabajo Fin de Máster, Universidad Politécnica de madrid, 2017): se trata de un estudio que tiene  similitudes en cuanto a que se realiza un estudio de la contaminación, se presentan las misma partículas y contaminantes, pero resulta curioso ver como varían las legislaciones al tratarse de un país de América Latina.
	\item \href{https://github.com/Fictizia/aireMAD}{\textit{AireMAD}} (Fictizia escuela, 2017): aplicación libre desarrollada por Fictizia (escuela de programación) para poder ver los datos de la calidad del Aire de Madrid en tiempo real. Actualmente no se encuentra en desarrollo pero tiene recursos interesante y una robusta documentación disponible  \href{https://github.com/Fictizia/aireMAD/wiki}{aquí}.
\end{itemize}


 

\capitulo{7}{Entorno experimental}

En esta sección se va a exponer de manuscrita el entorno con el que se ha realizado el estudio. 


\section{Esquema}

A modo de introducción se han seguido cuatro partes principales en el desarrollo técnico del proyecto en las que se han utilizado diferentes técnicas. Los siguientes esquemas (ver figuras \ref{diagramtechA} y \ref{diagramtechB}) describen de manera la línea a seguir en el proyecto.

   \imagen{diagramtechA}{\footnotesize{Detalle diagrama desarrollo (Partes 1 y 2). Fuente: Elaboración propia.}\label{diagramtechA}}
    \imagen{diagramtechB}{\footnotesize{Detalle diagrama desarrollo (Partes 3 y 4). Fuente: Elaboración propia.}\label{diagramtechB}}

\section{Fuente datos origen}

En el presente proyecto se han utilizado \textit{datasets} obtenidos principalmente a través de portales de datos abiertos, siendo esta la principal fuente de datos. Estos portales ponen a disposición pública múltiples sets de datos siguiendo una estrategia basada en cuatro aspectos fundamentales: la apertura de datos, la transparencia, la interacción y la participación de personas y/o empresas. 

El portal de Datos Abiertos del Ayuntamiento de Madrid, se dedica a promover el acceso a los datos del gobierno municipal y trata de impulsar el desarrollo de herramientas creativas para atraer y servir a la ciudadanía de Madrid. 
Este portal dispone de un amplio catálogo de datos que puede ser descargado a través de un acceso web para el público general (descarga ordinaria) y un \textit{API REST} con el que automatizar y programar el acceso y descarga de los diferentes \textit{datasets}. A través de este portal \cite{portal_datosabiertos_madrid}, es posible descargar un set de datos de calidad del aire con mediciones horarias y diarias desde el año 2001 hasta la actualidad. También existen datos en tiempo real.

Por lo tanto los tres tipos de ficheros que podemos descargar son los siguientes:

\begin{itemize}
	\item \textbf{Datos de calidad del aire horarios}: en este conjunto de datos se puede obtener la información recogida por las estaciones de control de calidad del aire, con los \textbf{datos diarios} por anualidades de 2001 a 2019. 
	\item \textbf{Datos de calidad del aire horarios:} en este conjunto de datos se puede obtener la información recogida por las estaciones de control de calidad del aire, con los \textbf{datos horarios} por anualidades de 2001 a 2019. Los datos horarios de las magnitudes corresponden a la media aritmética de los valores diezminutales que se registran cada hora.
	\item \textbf{Datos de calidad del aire en tiempo real:} En este conjunto de datos se puede obtener la información actualizada en tiempo real.
\end{itemize}

En el caso concreto del proyecto nos vamos a centrar en los dos primeros. En el caso de los \textbf{datos horarios}, éstos viene comprimidos en un archivo \textit{.zip} y cuando nos los descargamos aparecen divididos por meses. El nombre del archivo esta formado por las tres letras primeras, que corresponden al mes más un guión bajo seguido de dos números \textit{XX}, que identifican al año.

En caso de los \textbf{datos diarios} es similar pero nos proponen el formato de descarga directamente desde la web.

Los principales aspectos a destacar de la fuente de datos serían: 
\begin{itemize}
	\item Los ficheros están en texto plano, \textit{.csv} o \textit{.xml} y los campos no se encuentran delimitados.
	\item Caso\textbf{ datos diarios}: los datos se encuentran almacenados mensualmente y agrupados en ficheros comprimidos por cada año.
	\item Caso\textbf{ datos horarios}: cada registro contiene los 24 valores horarios de un día, 30 ó 31 filas contiguas corresponden a los valores de los días del mes, repitiéndose con cada magnitud (contaminante) de todas las estaciones que lo miden. Cada fichero contiene un mes de observaciones.
	\item Se dispone de mediciones de calidad del aire desde el año 2001 hasta la actualidad. 
	\item Cada campo tiene asignado un determinado número de dígitos. 
	\item Todos los campos contienen datos numéricos ya sean identificadores o medidas.
	\item Hay que precisar que las medidas válidas están marcadas con una V y las medidas no válidas están marcadas con una N. Tan solo serán validas las que contienen la V.
	\item Los datos anteriores a 2010 presentan algunos errores y los formatos están delimitados
\end{itemize}

La tabla 7.5 contiene los código de interpretación de los \textbf{datos horarios}. Tres apuntes importantes:

\begin{enumerate}
\item El campo punto de muestreo incluye el código de la estación completo (provincia, municipio y estación) más la magnitud y la técnica de muestreo.

\item H01 corresponde al dato de la 1 de la mañana de ese día, V01 es el código de validación, H02 al de las 2 de la mañana, V02 y así sucesivamente.

\item Únicamente son válidos los datos que llevan el código de validación “V".
\end{enumerate}


%DATOS HORARIOS
\begin{table}[H]
	\begin{center}
		\begin{tabular}{|c|c|}
			\hline
			\textbf{PROVINCIA}       & \textit{28}              \\ \hline
			\textbf{MUNICIPIO}       & \textit{79}              \\ \hline
			\textbf{ESTACION}        & \textit{4}               \\ \hline
			\textbf{MAGNITUD}        & \textit{1}               \\ \hline
			\textbf{PUNTO\_MUESTREO} & \textit{28079004\_1\_38} \\ \hline
			\textbf{ANO}             & \textit{2019}            \\ \hline
			\textbf{MES}             & \textit{1}               \\ \hline
			\textbf{DIA}             & \textit{1}               \\ \hline
			\textbf{H01}             & \textit{23}              \\ \hline
			\textbf{V01}             & \textit{V}               \\ \hline
			\textbf{H02}             & \textit{17}              \\ \hline
			\textbf{V02}             & \textit{V}               \\ \hline
			\textbf{[...]}             & \textit{[...]}               \\ \hline
		\end{tabular}
		\caption{Tabla intérprete valores horarios.}
	\end{center}
\label{TablaValHorarios}
\end{table}

La tabla 7.6 contiene los código de interpretación de los \textbf{datos diarios}. Tres apuntes importantes:

\begin{enumerate}
	\item El campo punto de muestreo incluye el código de la estación completo (provincia, municipio y estación) más la magnitud y la técnica de muestreo.
	
	\item D01 corresponde al dato del primer día del mes, D02 al del segundo día y así sucesivamente.
	
	\item Únicamente son válidos los datos que llevan el código de validación “V".
\end{enumerate}

%VALORES DIARIOS
\begin{table}[H]
	\begin{center}
		\begin{tabular}{|c|c|}
			\hline
			\textbf{PROVINCIA}       & \textit{28}              \\ \hline
			\textbf{MUNICIPIO}       & \textit{79}              \\ \hline
			\textbf{ESTACION}        & \textit{4}               \\ \hline
			\textbf{MAGNITUD}        & \textit{1}               \\ \hline
			\textbf{PUNTO\_MUESTREO} & \textit{28079004\_1\_38} \\ \hline
			\textbf{ANO}             & \textit{2019}            \\ \hline
			\textbf{MES}             & \textit{1}               \\ \hline
			\textbf{D01}             & \textit{18}              \\ \hline
			\textbf{V01}             & \textit{V}               \\ \hline
			\textbf{D02}             & \textit{20}              \\ \hline
			\textbf{V02}             & \textit{V}               \\ \hline
			\textbf{[...]}             & \textit{[...]}               \\ \hline
		\end{tabular}
	\caption{Tabla intérprete valores diarios.}
	\end{center}
\label{TablaValDiarios}
\end{table}


Como ya hemos nombrado anteriormente, el Sistema de Vigilancia está formado por 24 estaciones remotas automáticas que recogen la información básica para la vigilancia atmosférica. Poseen los analizadores necesarios para la medida correcta de los niveles de gases y de partículas. La localización de las estaciones de control también se encuentra para su libre descarga, este caso se permiten los ficheros \textit{.xls}, \textit{.csv} y \textit{.geo}. La tabla 7.7 contiene los datos para la correcta interpretación de los ficheros.


%TABLA ESTACIONES
\begin{table}[H]
			\begin{center}
	\begin{tabular}{|c|c|c|}
	
		\hline
		\textbf{Código estación} & \textbf{Lugar}       & \textbf{Comentarios}         \\ \hline
		\textit{28079001}        & Pº. Recoletos        & Baja.- 04/05/2009 (14:00 h.) \\ \hline
		\textit{28079002}        & Glta. de Carlos V    & Baja.- 04/12/2006 (11:00 h.) \\ \hline
		\textit{28079003}        & Pza. del Carmen      & * Código desde enero 2011    \\ \hline
		\textit{28079004}        & Pza. de España       &                              \\ \hline
		\textit{28079005}        & Barrio del Pilar     & * Código desde enero 2011    \\ \hline
		\textit{28079006}        & Pza. Dr. Marañón     & Baja.- 27/11/2009 (08:00 h.) \\ \hline
		\textit{28079007}        & Pza. M. de Salamanca & Baja.- 30/12/2009 (14:00 h.) \\ \hline
		\textit{28079008}        & Escuelas Aguirre     &                              \\ \hline
		\textit{28079009}        & Pza. Luca de Tena    & Baja.- 07/12/2009 (08:00 h.) \\ \hline
		\textit{28079010}        & Cuatro Caminos       & * Código desde enero 2011    \\ \hline
		\textit{28079011}        & Av. Ramón y Cajal    &                              \\ \hline
		\textit{28079012}        & Pza. Manuel Becerra  & Baja.- 30/12/2009 (14:00 h.) \\ \hline
		\textit{28079013}        & Vallecas             & * Código desde enero 2011    \\ \hline
		\textit{28079014}        & Pza. Fdez. Ladreda   & Baja.- 02/12/2009 (09:00 h.) \\ \hline
		\textit{28079015}        & Pza. Castilla        & Baja.- 17/10/2008 (11:00 h.) \\ \hline
		\textit{28079016}        & Arturo Soria         &                              \\ \hline
		\textit{28079017}        & Villaverde Alto      &                              \\ \hline
		\textit{28079018}        & C/ Farolillo         &                              \\ \hline
		\textit{28079019}        & Huerta Castañeda     & Baja.- 30/12/2009 (13:00 h.) \\ \hline
		\textit{28079020}        & Moratalaz            & * Código desde enero 2011    \\ \hline
		\textit{28079021}        & Pza. Cristo Rey      & Baja.- 04/12/2009 (14:00 h.) \\ \hline
		\textit{28079022}        & Pº. Pontones         & Baja.- 20/11/2009 (10:00 h.) \\ \hline
		\textit{28079023}        & Final C/ Alcalá      & Baja.- 30/12/2009 (14:00 h.) \\ 	\hline
		\textit{28079024} & Casa de Campo           &                                                                                                  \\ \hline
		\textit{28079025} & Santa Eugenia           & Baja.- 16/11/2009 (10:00 h.)                                                                     \\ \hline
		\textit{28079026} & Urb. Embajada (Barajas) & Baja.- 11/01/2010 (09:00 h.)                                                                     \\ \hline
		\textit{28079027} & Barajas                 &                                                                                                  \\ \hline
		\textit{28079047} & Méndez Álvaro           & Alta.- 21/12/2009 (00:00 h.)                                                                     \\ \hline
		\textit{28079048} & Pº. Castellana          & Alta.- 01/06/2010 (00:00 h.)                                                                     \\ \hline
		\textit{28079049} & Retiro                  & Alta.- 01/01/2010 (00:00 h.)                                                                     \\ \hline
		\textit{28079050} & Pza. Castilla           & Alta.- 08/02/2010 (00:00 h.)                                                                     \\ \hline
		\textit{28079054} & Ensanche Vallecas       & Alta.- 11/12/2009 (00:00 h.)                                                                     \\ \hline
		\textit{28079055} & Urb. Embajada (Barajas) & Alta.- 20/01/2010 (15:00 h.)                                                                     \\ \hline
		\textit{28079056} & Plaza Elíptica          & Alta.- 18/01/2010 (12:00 h.)                                                                     \\ \hline
		\textit{28079057} & Sanchinarro             & Alta.- 24/11/2009 (00:00 h.)                                                                     \\ \hline
		\textit{28079058} & El Pardo                & Alta.- 30/11/2009 (13:00 h.)                                                                     \\ \hline
		\textit{28079059} & Parque Juan Carlos I    & Alta.- 14/12/2009 (00:00 h.)                                                                     \\ \hline
		\textit{28079086} & Tres Olivos    & Alta.- 14/01/2010 (13:00 h.)                                                                    \\ \hline
	
\end{tabular}
	\caption{Tabla intérprete estaciones de control.}
	\end{center}	
\end{table}

Y para terminar los contaminantes, descritos en la tabla 7.8, también es necesario interpretarlos.

%TABLA CONTAMINANTES
\begin{table}[H]
		\begin{center}
	\begin{tabular}{|c|c|c|c|c|c|}
		\hline
		\multicolumn{2}{|c|}{\textbf{Magnitud}}     & \textbf{Abreviatura} & \textbf{Unidad} & \multicolumn{2}{c|}{\textbf{Técnica de medida}} \\ \hline
		01 & \textit{Dióxido de Azufre}             & \ce{SO2}                  & \textmugreek g/m3    & 38               & Fluorescencia ultravioleta         \\ \hline
		06 & \textit{Monóxido de Carbono}           & \ce{CO}                   & mg/m3   & 48               & Absorción infrarroja               \\ \hline
		08 & \textit{Dióxido de Nitrógeno}           & \ce{NO2}                   & \textmugreek g/m3           & 48         &  Id.               \\ \hline
		
		07 & \textit{Monóxido de Nitrógeno}         & \ce{NO}                   & \textmugreek g/m3           & 08         & Quimioluminiscencia                \\ \hline
		09 & \textit{Partículas \textless 2.5 \textmugreek m}   & \ce{PM2,5}                & \textmugreek g/m3           & 47         & Microbalanza                       \\ \hline
		10 & \textit{Partículas \textless 10 \textmugreek m}    & \ce{PM10}                 & \textmugreek g/m3           & 47         & Id.                                \\ \hline
		12 & \textit{Óxidos de Nitrógeno}           & \ce{NOx}                  & \textmugreek g/m3           & 08         & Quimioluminiscencia                \\ \hline
		20 & \textit{Tolueno}                       & \ce{TOL}                  & \textmugreek g/m3           & 59         & Gases                              \\ \hline
		30 & \textit{Benceno}                       & \ce{BEN}                  & \textmugreek g/m3           & 59         & Id.                                \\ \hline
		35 & \textit{Etilbenceno}                   & \ce{EBE}                  & \textmugreek g/m3           & 59         & Id.                                \\ \hline
		37 & \textit{Metaxileno}                    & \ce{MXY}                  & \textmugreek g/m3           & 59         & Id.                                \\ \hline
		38 & \textit{Paraxileno}                    & \ce{PXY}                  & \textmugreek g/m3           & 59         & Id.                                \\ \hline
		39 & \textit{Ortoxileno}                    & \ce{OXY}                  & \textmugreek g/m3           & 59         & Id.                                \\ \hline
		42 & \textit{Hidrocarburos totales(hexano)} & \ce{TCH}                  & mg/m3           & 02         & Ionización de llama                \\ \hline
		43 & \textit{Metano}                        & \ce{CH4}                  & mg/m3           & 02         & Id.                                \\ \hline
		44 & Hidrocarburos no metánicos             & \ce{NMHC}                & mg/m3           & 02         & Id.                                \\ \hline
	\end{tabular}
	\caption{Tabla ínterprete contaminantes.}
\end{center}	
\end{table}


\section{Exploración}
Tras conocer, estudiar a fondo los tipos de datos existentes y su correcta interpretación es tiempo de explorarlos. La fase de exploración se ha realizado de diferentes maneras y desde varias perspectivas con la ayuda de diversas herramientas \textit{open source}.

Esta fase es necesaria para realizar una investigación preliminar con el fin de entender mejor las características específicas de nuestros datos. En esta fase buscaremos correlaciones, tendencias y valores atípicos. Sin esta fase, no podríamos utilizar los datos de manera eficaz. Hemos hecho uso de diversas herramientas como es el caso de \href{http://vega.github.io/}{Voyager2} (figura \ref{voyager}).

    \imagen{voyager}{\footnotesize{Detalle \textit{vega-voyager}. Fuente: Elaboración propia.}\label{voyager}}
    
Datawrapper (ver figura \ref{datawrapper}) es otra herramienta que nos ayuda a analizar nuestros datos de una manera rápida. Es conveniente no coger una muestra demasiado grande ya que, ni es necesario, ni el software está correctamente optimizado para grandes volúmenes.
    
        \imagen{datawrapper}{\footnotesize{Detalle \textit{Datawrapper}. Fuente: Elaboración propia.}\label{datawrapper}}
 
A continuación, en las tablas 7.9 y 7.9, se expone la tipología de las variables en su conjunto. Aunque es posible que no se utilicen todas, sí es necesario introducir aquí todas las variables para establecer una relación con lo descrito en páginas anteriores.

    %DATOS HORARIOS
    \begin{table}[H]
    	\begin{center}
    		\begin{tabular}{|c|c|}
    			\hline
    			\textbf{PROVINCIA}       & \textit{integer}              \\ \hline
    			\textbf{MUNICIPIO}       & \textit{integer}              \\ \hline
    			\textbf{ESTACION}        & \textit{integer}               \\ \hline
    			\textbf{MAGNITUD}        & \textit{integer}               \\ \hline
    			\textbf{PUNTO\_MUESTREO} & \textit{integer} \\ \hline
    			\textbf{ANO}             & \textit{integer}            \\ \hline
    			\textbf{MES}             & \textit{integer}               \\ \hline
    			\textbf{DIA}             & \textit{integer}               \\ \hline
    			\textbf{H01}             & \textit{integer}              \\ \hline
    			\textbf{V01}             & \textit{float}               \\ \hline
    			\textbf{H02}             & \textit{integer}              \\ \hline
    			\textbf{V02}             & \textit{text}               \\ \hline
    			\textbf{[...]}             & \textit{[...]}               \\ \hline
    		\end{tabular}
    		\caption{Tabla tipo datos valores horarios.}
    	\end{center}
    \end{table}

%DATOS DIARIOS
\begin{table}[H]
	\begin{center}
		\begin{tabular}{|c|c|}
			\hline
			\textbf{PROVINCIA}       & \textit{integer}              \\ \hline
			\textbf{MUNICIPIO}       & \textit{integer}              \\ \hline
			\textbf{ESTACION}        & \textit{integer}               \\ \hline
			\textbf{MAGNITUD}        & \textit{integer}               \\ \hline
			\textbf{PUNTO\_MUESTREO} & \textit{integer} \\ \hline
			\textbf{ANO}             & \textit{integer}            \\ \hline
			\textbf{MES}             & \textit{integer}               \\ \hline
			\textbf{D01}             & \textit{integer}              \\ \hline
			\textbf{V01}             & \textit{text}               \\ \hline
			\textbf{D02}             & \textit{integer}              \\ \hline
			\textbf{V02}             & \textit{text}               \\ \hline
			\textbf{[...]}             & \textit{[...]}               \\ \hline
		\end{tabular}
		\caption{Tabla tipo datos valores diarios.}
	\end{center}
\end{table}

Para conocer más a fondo nuestros datos algunas estadísticas básicas de este tipo que podemos calcular para nuestro conjunto de datos son la media, mediana, el rango y la desviación estándar. Son muy interesantes para este estudio los promedios y la media de los gases contaminantes además de los máximos. En numerosas ocasiones organismos como la OMS \cite{oms_1} dictaminan episodios de contaminación en función de la media cada\textit{ x} período horario. Cada contaminante tiene sus medidas y cada organismo dictamina su correcta interpretación.

 
\section{Transformación de datos}
Tras explorar y entender los datos es hora de transformarlos de acuerdo a nuestros objetivos finales. Es cierto que se han utilizado diferentes lenguajes de programación pero en todos se han seguido patrones comunes a la hora de transformar los datos.

Algunas características comunes de la transformación de los datos han sido el renombramiento de columnas, el establecer el tipo de dato adecuado por columnas para realizar operaciones y evitar problemas a la hora de realizar operaciones, el intercambiar las columnas de validación para establecer los datos que son correctos o el agrupar por contenido para obtener una determinada gráfica. Además se han limpiado las columnas no válidas, se han desechado las que no iban a usarse y se han agrupado las columnas de tipo fecha. 

Por otro lado, en algunas ocasiones los cálculos de promedios, medias o recuentos se han realizado directamente sobre \textit{PowerBi.} Los niveles marcados por la OMS en algunas ocasiones son fijos por lo que también se han añadido de manera manual cuando ha sido necesario. 

Vamos a poner un ejemplo del procedimiento a seguir para \href{https://tinybird.co//}{Tinybird.co}. Lo primero que hay que hacer es cargar los datos. Como hemos nombrado se trata de una aplicación en beta por lo tanto el entorno gráfico no está disponible todavía para subir cantidades  grandes de datos por lo que se ha creado un script en bash que a través de un bucle sube todos los archivos. 

\lstset{breaklines=true, basicstyle=\footnotesize}
\begin{lstlisting}[frame=single]
#!/bin/bash

#Personal Token
TOKEN=p.eyXXXXXXXXXXXXXXXXXXXXXX
#URL
BASE_URL=/Users/xxxx/Documents//TFM_MUINBDES/Code/data

#Load months
for month in ene feb mar abr may jun jul ago sep oct nov dic
do
curl -F "csv=@${month}_mo18.csv" \
-H "Authorization: Bearer $TOKEN" \
-X POST "https://XXXXXXXXXXX/v0/datasources?name=datos_2018&mode=append"
done
\end{lstlisting}

Una vez tenemos los datos cargados (ver figura \ref{tiny1}) es el momento de transformar los datos de acuerdo a nuestro estudio particular. 

\imagen{tinybird1}{\footnotesize{Datos cargados en \textit{tinybird}. Fuente: Elaboración propia.}\label{tiny1}}

En este caso queremos observar si existen patrones anuales en los datos de la cantidad de \ce{NO2} para después comprobar la misma hipótesis de manera semanal y por horas. Lo primero que hacemos es seleccionar de nuestro conjunto de datos lo que nos hace falta, además transformamos la fecha (ya que recordamos que la fecha viene en nuestro fichero de datos original separada en tres columnas diferentes) para una mejor visualización (ver figura \ref{tiny2}).

\imagen{tinybird2}{\footnotesize{Detalle \textit{tinybird}. Fuente: Elaboración propia.}\label{tiny2}}

Después escogemos los valores válidos y también la magnitud que corresponde al contaminante del cuál queremos hacer el análisis (ver figura \ref{tiny3})

\imagen{tinybird3}{\footnotesize{Detalle \textit{tinybird}. Fuente: Elaboración propia.}\label{tiny3}}

Para terminar agrupamos y ordenamos por lo que nos hace falta para nuestro análisis (ver figura \ref{tiny4}) y ya podemos exportar el archivo para proceder a su visualización.

\imagen{tinybird4}{\footnotesize{Detalle \textit{tinybird}. Fuente: Elaboración propia.}\label{tiny4}}

 En el caso de python, cuyo archivo final se ha utilizado para la visualización interactiva, el proceso ha sido muy similar con la peculiaridad de que se han añadido los datos de las estaciones de medida al fichero final. En este caso, como ya he nombrado, se han escogido los datos de 2014 a 2018. Se han seleccionado los datos directamente desde la web para evitar inconsistencias. Hemos tenido que lidiar con los archivos en \textit{.txt}, ya que algunos años solo aparecían en este formato. Hemos hecho un \textit{merge} de los datos de las estaciones junto con los datos de los gases para establecer por estación los tipos de gases medidos. De esta manera se ha logrado enriquecer la visualización añadiendo correspondencia por zonas y estaciones.
 
 
 \begin{itemize}
 	\item Zona 1
 	 \begin{itemize}
 		\item Mendez Álvaro
 		\item Arganzuela
 		\item Carlos V
 		\item Parque del Retiro
 		\item Escuelas Aguirre
 		\item Plaza del Carmen
 		\item Paseo de Recoletos
 		\item Plaza de España
 		\item Manuel Becerra
 		\item Barrio Salamanca
 		\item Castellana
 		\item Gregorio Marañon
 		\item Cuatro Caminos
 		\item Ciudad Universitaria
 		\item Avenida Ramón y Cajal
 		\item Plaza Castilla
 		\item Barrio del Pilar
 	\end{itemize}
 	\item Zona 2
 	 \begin{itemize}
 		\item Moratalaz
 		\item Vallecas
 		\item Santa Eugenia
 		\item Ensanche de Vallecas
 	\end{itemize}
 	\item Zona 3
 	 \begin{itemize}
 		\item Arturo Soria
 		\item Campo de las Naciones
 		\item Juan Carlos I
 		\item Urbanización Embajada
 		\item Barajas Pueblo
 		\item San Chinarro
 		\item Tres Olivos Plaza
 	\end{itemize}
 	\item Zona 4
 	 \begin{itemize}
 		\item El Pardo
 		\item Casa de Campo
 		\item Alto de Extremadura
 	\end{itemize}
 	\item Zona 5
 	 \begin{itemize}
 		\item Farolillo
 		\item Cuatro Vientos
 		\item Plaza Fernández Ladreda (Elíptica)
 		\item Villaverde
 	\end{itemize}
 \end{itemize}

 

\section{Visualización de datos}
Una vez que tenemos los datos transformados de manera adecuada ya podemos concentrarnos en la parte visual. En este trabajo la parte visual está dividia en varias partes, por un lado en todas las gráficas que hemos empleado para explorar los datos; segundo, en una visualización final a modo de dashboard realizada con \textit{PowerBi};  y en una pequeña web resumen de resultados.


Siguiendo con el ejemplo nombrado en el apartado anterior y cargando los datos en \textit{vega} podemos obtener la visualización de acuerdo a nuestro objetivo. En este caso hemos calculado la media de \ce{NO2} por mes (en naranja) junto con los máximos (en azul)  del año 2019, concretamente de enero a julio. En la siguiente figura \ref{NO22019} podemos ver el gráfico de líneas.

\imagen{NO22019}{\footnotesize{Detalle \ce{NO2} 2019. Fuente: Elaboración propia.}\label{NO22019}}


También lo mismo pero por día de la semana (ver figura \ref{ejemploJulio}), en este caso para el mes de julio y con un gráfico de barras.

\imagen{ejemploJulio}{\footnotesize{Detalle \ce{NO2} 2019 (julio). Fuente: Elaboración propia.}\label{ejemploJulio}}


Y lo mismo pero por horas, en este caso para el 1 de julio (ver figura \ref{ejemplo01julio}) y con un gráfico de puntos.
\imagen{ejemplo01julio}{\footnotesize{Detalle \ce{NO2} 2019 (1 julio). Fuente: Elaboración propia.}\label{ejemplo01julio}}

Con estas tres figuras anteriores ya comenzamos a ver patrones en los datos de los cuales hablaremos en la parte de resultados.

La parte de la visualización en \textbf{PowerBi} es más interactiva y refleja una parte de los datos, 2014 a 2018, desde el punto de vista de la OMS y los gases que esta organización refleja como más perjudiciales para la salud.  Una vez transformados los datos y con el fichero limpio se cargan los datos. He realizado distintas páginas dentro de \textit{PowerBi }en las que podemos ver diferentes tipos de gráficos.

\begin{itemize}
	\item \textbf{Página 1}: Portada del trabajo
	\item \textbf{Página 2}: Resumen de todos los contaminantes por zonas y años. De este modo se puede ver la variación a lo largo del tiempo de todos los contaminantes principales.
	\item \textbf{Página 3} (ver figura \ref{PowerBi2}): Esta página-resumen es la misma que la anterior pero en esta ocasión nos centramos tan solo en  \ce{NO2} y se expone la variabilidad por días de la semana pudiendo ver de manera rápida los días más afectados.
	\imagen{pagina2PowerBi.png}{\footnotesize{Detalle gráfico por día de la semana. Fuente: Elaboración propia.}\label{PowerBi2}}

	\item \textbf{Página 4}: \textit{límites anuales (OMS)},	se exponen aquí la media de las partículas por millón de \ce{PM2,5} y \ce{PM10} , además de  \ce{NO2}.  La línea azul refleja el nivel recomendado por la OMS. Es posible filtar por zonas y ver en las gráficas de barras donde se superan los niveles. Se ve que se respetan los niveles en casi todos los casos y en otros vemos que se supera lo establecido por la OMS.
	
	\item \textbf{Página 5} (ver figura \ref{pagina5}):  \textit{límites diarios (OMS)}, en este caso volvemos a escoger las partículas por millón de \ce{PM2,5} y \ce{PM10} , además de \ce{NO2}. Es importante conocer aquí que dependiendo del contaminante se establecen unos límites distintos por la OMS. En este caso se establecen los límites por día. La línea azul refleja ese límite y los gráficos de color naraja la media diaria para esos valores seleccionados. Es posible también filtrar por zonas y estaciones. 
	
	\imagen{pagina5}{\footnotesize{Detalle límite diario OMS. Fuente: Elaboración propia.}\label{pagina5}}
	
	\item \textbf{Página 6}: nos centramos en los dos contaminantes más criticos según la OMS, el ozono (\ce{O3}) y el dióxido de nitrógeno (\ce{NO2}). La OMS recomienda un máximo para la media cada ocho horas. Se reflejan en los gráficos ambos contaminantes respectivamente en el primer y segundo gráfico. Como se puede ver superamos con creces los niveles de \ce{O3}.
	
	\imagen{pagina6}{\footnotesize{Detalles \ce{O3} y \ce{NO2}  . Fuente: Elaboración propia.}\label{pagina6}}
	
	\item \textbf{Página 7}: El procolo anticontaminación es superado en numerosas ocasiones por los diferentes niveles que son los que ponen en alerta al Ayuntamiento. En este informe podemos ver el número de veces que se supera por zonas. El objetivo es poder filtrar con detalle por zona, estación y período.  
	
	\item \textbf{Página 8}: La media de \ce{NO2} es la que se utiliza para ver si se han superado esos niveles de los que hablábamos. En esta ocasión podemos ver dos gráficos de líneas que nos sirven para comparar entre estaciones y por período/fecha si se han superado los niveles.
	
	\item \textbf{Página 9} (ver figura \ref{MapaNO2}): Es posible ver aquí otra interpretación de la página anterior. En esta ocasión se representa la  acumulación de \ce{NO2} por zonas y horas de manera visual. Concretamente se hace uso de un mapa utilizando un \textit{plugin} de mapbox. 
	
	En color amarillo los que ya han superado ese 180 (preaviso), en color rojo los que han superado el 200 (aviso) y en color verde los que no llegan a esos límites. La acumulación cambia y, tal y como recogen los informes del ayuntamiento, se produce a medida que el tráfico aumenta por lo que tiene un efecto tardío.
	
	\imagen{MapaNO2}{\footnotesize{Detalle mapa \ce{NO2} (\textit{MapBox}). Fuente: Elaboración propia.}\label{MapaNO2}}
	
	\item \textbf{Página 10} (ver figura \ref{NivelProtocolo}): para terminar un mapa de calor que indica el nivel del protocolo para los preavisos (mayor de 180). Sigue la línea del mapa anterior y es posible filtar por zona, hora del día y período.
	
	\imagen{NivelProtocolo}{\footnotesize{Detalle nivel protocolo (\textit{MapBox}). Fuente: Elaboración propia.}\label{NivelProtocolo}}
	
	\item \textbf{Páginas extra}: también se ha experimentado con otro tipo de gráficos (en este caso siempre con niveles de \ce{NO2}) en los que se puede ver de manera simple, cual es la zona más afectada por tamaño de las palabras, figura \ref{nubePalabras}, en un gráfico de palabras; o dependiendo de la estación/año, figura \ref{powerBitamanos}, en un gráfico de rectángulos. También es posible ver cuáles son los días por semana que más cantidad promedio tienen, figura \ref{powerBidias}, y además si se han superado los avisos. Las posibilidades son infinitas.
	
		\imagen{nubePalabras}{\footnotesize{Ejemplo gráfico nube de palabras.}\label{nubePalabras}}
	
	
	\imagen{powerBitamanos}{\footnotesize{Ejemplo gráfico por tamaño/año. Fuente: Elaboración propia.}\label{powerBitamanos}}
	

	\imagen{powerBidias}{\footnotesize{Ejemplo gráfico por días y zonas. Fuente: Elaboración propia.}\label{powerBidias}}
	
\end{itemize}



\capitulo{8}{Resultados}

En está penúltima sección se exponen las conclusiones derivadas del trabajo. 



\section{Preámbulo}
Los resultados son el remate final a este estudio, sin embargo, quiero hacer aquí una pequeña aclaración. Los efectos de este estudio realizado por otro conjunto de personas y partiendo de los mismos datos pueden ser diferentes dado que el enfoque que se le da es distinto. Por ejemplo yo me he enfocado más en un contaminante, \ce{NO2}, y en éste comparado a los niveles de la OMS \cite{oms_2}. Si bien es cierto que se ha encontrado patrones similares con otros trabajos (ver más en la sección 6) no era este el objetivo final. 

En la figura \ref{diagramtech} se puede ver el conjunto de desarrollo completo que se ha seguido para llegar a las visualizaciones.

  \imagen{diagramtech}{\footnotesize{Diagrama desarrollo. Fuente: Elaboración propia.}\label{diagramtech}}
  
  Los resultados que se exponen en el siguiente apartado son los que se han obtenido de observar, analizar y estudiar esas visualizaciones para dar respuesta a las preguntas analíticas. Una vez con los datos limpios los gráficos o resultados a obtener son tantos como objetivos tengamos en mente.
  

\section{Resultados generales}

Después de un análisis exhaustivo de los datos para los últimos cinco años, podemos concluir con que Madrid tiene un problema de contaminación severo. No solo con los datos en la mano se ratifica esta postura sino que desde la unión Europea 

Además según la OMS la práctica totalidad de la población de España viene respirando aire contaminado, con foco principal en Madrid como punto elevado de concentración de gases tóxicos. Con los datos delante se aprecia una clara mejoría en los datos y en el año 2018 se mejoran los datos respecto del año anterior, 2017.

Una de las primeras dudas que nos pueden surgir a la hora de hacer este tipo de estudios es por qué elegir un tipo de gas u otros. Ya no solo por que la OMS sabe de lo que habla y recoge en su legislación los gases más contaminantes para el ser humano sino basta ver la gráfica promedio de contaminantes (ver figura \ref{promTODOS}) para darse cuenta de los gases que mayores niveles tienen en los medidores son los más perjudiciales.

\imagen{promTODOS}{\footnotesize{Detalle gases tóxicos (2014-2018). Fuente: Elaboración propia.}\label{promTODOS}}

De acuerdo a ese tipo de gráficas promedio por contaminante el incremento de determinados tipo de contaminantes como las \ce{PM2,5} (Partículas Menores de 2,5 \textmugreek m), \ce{PM10} (Partículas Menores de 10 \textmugreek m), \ce{NO} (monóxido de nitrógeno) y \ce{NO2} (dióxido de nitrógeno) se viene relacionando desde hace tiempo con el aumento de ingresos hospitalarios por afecciones respiratorias y cardiovasculares por lo que tiene sentido el estudio de estas particulas.

Más concretamente nos hemos centrado en \ce{NO2}. Algunas apreciaciones interesantes que se han obtenido son por ejemplo que influye la \textbf{estacionalidad}. Los valores se disparan en los períodos comprendidos entre mediados de octubre y mediados de enero. Ver figura \ref{estacionalidad} con la media anual del año 2018 para ver claramente los patrones. Los meses con los menores niveles son junio, julio y agosto. 

\imagen{estacionalidad}{\footnotesize{Media \ce{NO2} en 2018. Fuente: Elaboración propia.}\label{estacionalidad}}

De la misma manera se observan patrones influyentes en el \textbf{día de la semana} y la \textbf{hora del día} . Se diferencian los días laborales de los festivos.  En días laborables existen dos zonas horarias de mayor contaminación correspondiendo éstas a las horas pico de tráfico por la mañana y por la tarde en torno a las 20-21 horas.  Se ve claramente en los días laborables como se producen dos picos claros que representar la entrada y salida del trabajo (ver figura \ref{promedioHora}).  En el caso del fin de semana, la zona más relevante se corresponde con las horas nocturnas, observándose un incremento progresivo desde las 19 horas hasta 24 horas. 

	
\imagen{promedioHora}{\footnotesize{Detalle promedio \ce{NO2} por hora. Fuente: Elaboración propia.}\label{promedioHora}}


La zona del \textbf{interior de la M30} es sin lugar a dudas donde se dan los valores más elevados de \ce{NO2}. También es el lugar dónde más afectados se ven los vehículos a partir del escenario de contaminación número 2. El mayor número de estaciones (más de un 40\%) se concentra en esta zona. De la misma manera el mayor número de valores que superan los 180 \textmugreek g/m3 también se sitúan en esta zona.

En la figura \ref{estaciones} es posible ver por orden de mayor a menor las estaciones dónde se producen los mayores niveles. Las líneas representan la cantidad total de preavisos (mayor 180) en color negro y en color rojo los avisos (mayor 200).

\imagen{estaciones}{\footnotesize{Detalle mayor concentración de \ce{NO2} por estación. Fuente: Elaboración propia.}\label{estaciones}}


Durante los años estudiados se producen grandes amplitudes en los datos aunque se focalizan por zonas como hemos dicho, siendo \textbf{dentro de la M30} donde se dan las mayores concentraciones. Cabe destacar que dentro de esta zona los valores menores se corresponden a la estación de Retiro. Esto podría confirmar que la masa forestal del Parque del Retiro disminuye sensiblemente los valores de contaminación.

Dentro de ésta zona podemos considerar que para todas las estaciones, salvo para la de Retiro, los valores atípicos son aquellos que están ligeramente por encima de 100-110  \textmugreek  g/m3.

Algo que hemos notado también es que el ayuntamiento de Madrid no aplica los protocolos siempre. Desde la página web del ayuntamiento se especifican cuando se aplican los protocolos \cite{mambiente} sin embargo en los datos podemos notar avisos durante la campaña de navidad de 2015 (ver figura \ref{navidad2015}, donde se aprecia claramente que se superan los niveles establecidos por OMS) en los cuales no se han activado los protocolos a pesar de tener avisos. Es posible que esas fechas en las que se produce un mayor número de desplazamientos el ayuntamiento tenga, quizás, algo de mano ancha. 

\imagen{navidad2015}{\footnotesize{Detalle período navidad 2015. Fuente: Elaboración propia.}\label{navidad2015}}

Se podría pensar que activando los protocolos se disminuiría la concentración de \ce{NO2}  (disminución del tráfico) para los días inmediatamente siguientes a la activación. Sin embargo, esto no se aprecia en picos altos durante el invierno. Por lo que se puede tener como hipótesis que parece que no basta el tráfico para la disminución de este gas. Puede deberse a qué durante el invierno un \% de emisiones de este gas provenga del uso de calefactores. Hay que tener en cuenta que, en el centro de las ciudades, que son las que se ven afectadas por los cortes de tráfico, la mayor parte de los edificios cuentan con sistemas de calefacción basadas en la combustión de carbón o gasóleo.

\section{Preguntas analíticas}

Algunas de las preguntas analíticas, marcadas al inicio del proyecto, han sido las siguientes:

\begin{itemize}
	\item \textit{¿La contaminación es un problema real en la ciudad de Madrid?} \textit{¿Supone un grave riesgo para la población?} \textit{¿Son las partículas de Madrid contaminantes de alto riesgo para la salud?}
	
	La principal respuesta a esta pregunta depende de los valores que se toman como referencia. Los valores que dictamina la OMS \cite{oms_2} son mucho más estrictos que los límites legales que establecen las instituciones y por lo tanto más acordes con la protección de la salud. De acuerdo a esos valores claramente Madrid tiene un serio riesgo de contaminación. 
	
	De acuerdo a los límites legales también dado que se ven superados en numerosas ocasiones por lo que los problemas para la salud existen. Además existen estudios \cite{ambiente_contaminacion_estudio} que ratifican el empeoramiento de la salud como causa directa de respirar este tipo de aires contaminados.
	
	Tal y como hemos estudiado en la ciudad de Madrid  los contaminantes que más incidencia presentan por regularidad y niveles son el dióxido de nitrógeno, el ozono  y las partículas \ce{PM10} y \ce{ PM2,5}. Los peores registros se lo llevan las estaciones que se encuentran dentro de la M30, por lo tanto las personas mas expuestas son las que se encuentran en esta zona pero por supuesto no las únicas.
	
	Madrid es de las pocas ciudades europeas en las que se sigue registrando el incumplimiento del valor límite horario de \ce{NO2}. Además desde Europa nos sigue llegando multas \cite{bruselas_} por incumplimiento del reglamento por lo que realmente parece que es un problema real en el que no se están tomando las medidas adecuadas. Si bien es cierto que no presenta niveles extremos como pueden darse en ciudades asiáticas, de cara al futuro convendría extremar las precauciones.
	
	Para terminar con la pregunta, y algo en relación al tema aunque fuera de la pregunta, según un último estudio \cite{calderon_garciduenas_combustion_2019} el corazón de las personas jóvenes que viven en grandes ciudades contienen miles de millones de partículas contaminantes. Estudio de la Universidad de Lancaster (Reino Unido), y que podría suponer la primera demostración de la ya conocida relación entre la calidad del aire y las enfermedades relacionadas con este órgano. El estudio se centra en diferentes partículas, como por ejemplo partículas derivadas de la exposición prolongada a \ce{NO2} u \ce{O3}, pero está desarrollado en la ciudad de México. Sería interesante ver un estudio de este tipo en otras ciudades Europeas, como Madrid.
	
	\item \textit{Los protocolos de anticontaminación se basan en los niveles de los óxidos de nitrógeno. Estos niveles suben por la noche por lo que siempre se espera a última hora para activar dichos protocolos. ¿Es una afirmación correcta?}
	
	Es una afirmación correcta con los datos en la mano. Pero vamos a analizarla un poco más. Las limitaciones, de diversas índoles dependiendo del escenario, comienzan a aplicarse partir de las 6:00 horas del día de entrada en vigor y termina cuando cesa el episodio de contaminación. Cuando este episodio cesa significa que una determinada estación no establece un preaviso, aviso o alerta pero eso no significa que la contaminación haya desaparecido. Normalmente 
	

	
	El problema que hemos detectado es que los protocolos de contaminación tan solo se aplican a determinadas zonas de Madrid. O al menos así se recogen en los informes públicos. Es decir, existen estaciones cuyas lecturas no activan, directamente, protocolos para restringir el tráfico u otras medidas. Los motivos son claros: al encontrarse fuera de la M-30, muchas estaciones que sobrepasan los límites establecidos quedan excluidos de los protocolos. Las alarmas pueden saltar por episodios de contaminación durante horas
	
	 Si nos paremos a pensar en factores externos como por ejemplo el clima, se dan casos en los que el centro de Madrid registra niveles seguros sin embargo todas las estaciones que lo rodean presentan preavisoscon a tasas insalubres sin que salten las alarmas.
	
	\item \textit{¿Existen contaminantes que puedan estar asociados a temporalidad?}
	
	Sin duda la respuesta es sí, un ejemplo muy claro es el del \ce{NO2}. Los efectos de este gas son producidos en su mayoría por la combustión. El \% mayor proviene de los coches pero durante el invierno existen multitud de hogares en la capital que todavía funcionan con viejas estufas de combustión por lo que los niveles de contaminación de las mediciones se ven afectado también por este tipo de episodios.
	
	\item \textit{¿Consideras que las últimas medidas aplicadas por el Ayuntamiento de Madrid durante los últimos años han mejorado la calidad del aire? ¿Qué contaminantes? ¿Alguna zona/estación en concreto? ¿Tiene sentido Madrid Central?}
	
	Los contaminantes más problemáticos en España durante los años se repiten de manera constante y son las	partículas en suspensión (\ce{PM10} y \ce{PM2,5}), el dióxido de nitrógeno (\ce{NO2}), el ozono troposférico (\ce{O3}) y el dióxido de azufre (\ce{SO2}). 
	
	De los últimos 6 años los peores registros se los llevan sin lugar a dudas 2017 pero en general se tienen en cuenta las medidas. Sin lugar a dudas la más polémica del pasado año ha sido Madrid Central. Si no estamos equivocados, el ayuntamiento de Madrid fue la primera ciudad en activar un protocolo \cite{protocolo} de estas características en nuestro país. 

	El 1 de julio de 2019 la corporación municipal de la ciudad cesó el conocido plan de la legislatura anterior para prevenir la contaminación en el centro de la ciudad. El nuevo ayuntamiento permite ahora entrar a cualquier automóvil para circular a sus anchas por Madrid Central, sin ningún tipo de
	limitación ni consecuencias. 
		
		Los datos de contaminación por \ce{NO2} registrados durante el mes de junio de 2019 en Madrid, muestran dos fases bien diferenciadas: una fase inicial en la que los valores de \ce{NO2} alcanzados fueron muy similares a los correspondientes a los meses de abril y mayo, es
		decir, con registros históricamente bajos, tanto en la estación ubicada dentro del perímetro de Madrid Central (recordemos es \textit{Plaza del Carmen}), como en el valor medio de toda la red; y una segunda fase, coincidiendo con el anuncio de la
		suspensión de las multas en Madrid Central a partir del 1 de julio, en la que se observa un rápido empeoramiento de los valores de \ce{NO2} registrados en la estación Plaza del Carmen, que se hace muy intensa en la última semana
		del mes. 
		
		En la figura \ref{junjul2019} tenemos exactamente las gráficas del mes de \textbf{Junio}, en la parte izquierda, con Madrid central aún en vigor; y del mes de \textbf{Julio}, en la parte derecha y con el fin de Madrid central. Se aprecian claramente los picos de los que hemos hablado en el párrafo anterior. En azul se aprencian los máximos y en naranja la media diaria de \ce{NO2} en cada mes.
		
		\imagen{junjul2019}{\footnotesize{Detalle \ce{NO2} Junio/Julio 2019. Fuente: Elaboración propia.}\label{junjul2019}}
		
		En la figura \ref{plazacarmen} se puede ver la evolución en años anteriores. Recordemos que Madrid Central entra en vigor a finales de 2018. Si comparamos la media de toda la red con respecto de la estación que se encuentra dentro de esta delimitación los valores de ésta última siempre han sido superiores hasta la entrada en vigor de las nuevas medidas. Éstas nuevas medidas se empiezan a notar a partir de enero de 2019.
		
		\imagen{plazacarmen}{\footnotesize{Detalle estación \textit{Plaza del carmen}. Fuente: Elaboración propia.}\label{plazacarmen}}
		
		Si por ejemplo comparamos entre \textit{Junio 2018} y \textit{Junio 2019}. En 2019 se registró 26 \textmugreek g/m3, 12 menos (reducción del 32 \%) que el registrado en 2018 (38 \textmugreek g/m3).
		
		 Realmente todavía es pronto para saber las consecuencias reales del cese de la prohibición dado que con un solo mes de comparativa, tan solo julio 2019 está disponible por el momento, resulta apresurado realizar estimaciones con acierto. También hay que tener en cuenta que nos encontramos en los meses de menores tasas por lo que será necesario esperar a los meses con valores más altos para determinar el impacto.
		 
		  Todavía ni los propios madrileños tienen claro si se aplican o no sanciones dentro de la zona cero ya que el juzgado de lo Contencioso-Administrativo número 7 de Madrid emitió una nueva resolución \cite{auto_judicial} que ratificaba la reactivación de multas en Madrid Central al dar el visto bueno a la medida cautelar solicitada por Greenpeace. Lo que sí resulta verídico es que desde Europa se plantean las sanciones administrativas \cite{20minutos_europa_2019} pertinentes en caso de no tomar medidas ante la calidad del aire de la ciudad de Madrid. 
		 
		 Lo que sí parecen los datos constatar es la mejora de la estación respecto de años anteriores durante el tiempo que Madrid Central seguía en vigor.
		
			

	
	\item \textit{¿Qué área es la más contaminada?}
	
	La área más contaminada es la del interior de la M30, y dentro de ella \textit{Escuelas Aguirre} tiene tendencia a marcar los mayores valores. Las estaciones de Pza. Fernández Ladreda (Suroeste) y las estaciones de Escuelas Aguirre (Interior M30) son las estaciones que suelen marcan los máximos. En relación a esta última podemos ver en la figura \ref{EscuelasAguirre} como claramente se superan los límites marcados por la OMS \cite{oms_2} en cuanto a los contaminantes partículas en suspensión (\ce{PM10} y \ce{PM2,5}) y el dióxido de nitrógeno (\ce{NO2}).
	
	\imagen{EscuelasAguirre}{\footnotesize{Detalle estación \textit{Escuelas aguirre}. Fuente: Elaboración propia.}\label{EscuelasAguirre}}
	
	Las diferencias entre el interior y el exterior de la M-30 son notables. En enero se captan los niveles más altos de todo el año y en los meses de verano, en los que la capital reduce su población debido a las vacaciones, estos niveles son reducidos. 
	
	
	\item \textit{¿Qué está ocurriendo en el año en curso 2019?} \textit{¿Demuestran los datos que son útiles las medidas implantadas?}
	
	El año en curso sigue la tendencia de mejora de los anteriores aunque se espera ver como va a afectar la medida del cese de Madrid Central. En la figura \ref{NO220192} se muestra la media de \ce{NO2} (en color naranja) durante los meses de enero a julio. La línea en color azul representa los máximos registrados.
	
	De lo que llevamos de año los meses de enero y febrero vienen tienendo valores similares a los de años anteriores en los mismos meses. En enero solo 3 estaciones redujeron la media respecto a otros años. Durante febrero la estación de Plaza del Carmen, la única ubicada dentro del área de Madrid Central, fue por primera vez en la que más disminuyó la contaminación respecto de otros años.
	
	\imagen{NO22019}{\footnotesize{Detalle \ce{NO2} 2019. Fuente: Elaboración propia.}\label{NO220192}}
	
	Tal y como es posible ver en la gráfica \ref{NO220192} en Marzo comienza un punto de inflexión. El tiempo comienza a variar ligeramente en esta época del año y es también el período (16 de marzo) donde se comienzan a imponer sanciones a los conductores que no respeten los protocolos. 
	
	Durante el mes de abril, en todas las estaciones salvo en dos (Juan Carlos I y Escuelas Aguirre), los niveles de contaminación disminuyeron respecto a años anteriores.
	
	En el mes de mayo todas las estaciones de medición acumularon resultados positivos; es decir, por primera vez desde la entrada en vigor de Madrid Central la ciudad tenía una calidad de aire por encima de la media del periodo de los cinco años anteriores.
\end{itemize}



	
\capitulo{9}{Conclusiones y Líneas de trabajo futuras}

En está última sección se exponen las conclusiones finales derivadas del trabajo. De la misma manera abordaremos las línea de trabajo futuro que considero que se pueden o deberían realizarse.


\section{Conclusiones proyecto}\label{conclusiones_proyecto}
Las conclusiones alcanzadas.
\begin{itemize}
	\item Datos de un entorno real estudiados, analizados, preprocesados y limpiados.
	\item \textit{Dashboard} final funcionando
	\item Página web con a modo de consulta  funcionado
	\item Patrones encontrados en los datos
	\item Preguntas analíticas planteadas respondidas
\end{itemize}

\section{Conclusiones personales}\label{conclusiones_personales}
Me gustaría exponer también unas conclusiones personales a modo de lista sobre lo que este trabajo ha supuesto para mi.  

\begin{itemize}
	\item Sin duda 
\end{itemize}

\section{Líneas de trabajo futuras}\label{lineas_futuras}
Es importante remarcar las línea de futuro que el proyecto pretende seguir.

\begin{itemize}
	\item Refactorizar 
\end{itemize}

%\renewcommand\chaptername{Anexo}
%\renewcommand\thechapter{\Roman{chapter}}
%\setcounter{chapter}{0}

% Añadir entrada en el índice: Anexos
\appendix
\addcontentsline{toc}{part}{Apéndices}
\part*{Apéndices}

%%Bibliografía
\bibliographystyle{plainnat}
\bibliography{bibliografia}



\apendice{Plan de Proyecto Software}

\section{Introducción}\label{introduccion-plan}
En este capitulo se detalla la planificación del proyecto. 

La planificación temporal del proyecto, o lo que es lo mismo, la elaboración del calendario o programa de tiempos, consiste en una representación gráfica de todas las actividades del proyecto necesarias para producir el resultado final que se desea.

Se han utilizado metodologías ágiles para el desarrollo del proyecto y de este modo, se ha realizado un desarrollo dividido en iteraciones. Terminada una iteración empezaba la siguiente y se agregaban a las tareas planeadas las que no habían sido completado de la iteracción precedente. Las iteraciones del proyecto, los sprints, estaban pensadas para durar unos \textbf{quince días} aproximadamente. No obstante, hay alguna excepción en la que la iteración duró más tiempo, o menos. También existe alguna demora entre algún sprint debido a que tenía demasiada carga de las asignaturas o demasiada carga laboral.

\section{Planificación temporal}\label{planificacion-temporal}
Desde el comienzo del proyecto se planteó utilizar una metodología ágil como
\emph{Scrum} para la gestión del proyecto. Aunque no se ha seguido al 100\% la
metodología al tratarse de un proyecto académico, sí que se ha aplicado
en líneas generales una filosofía ágil y metódica. La diferencia fundamental radica en que esta metodología esta pensada para equipos y no para individuos.

\subsection{Desarrollo con Scrum}

Scrum es un marco de trabajo que define un conjunto de prácticas y roles, y que puede tomarse como punto de partida para definir el proceso de desarrollo que se ejecutará durante el proyecto. Indicar aquí que se ha utlizado terminología en inglés dado que por convenio se siguen estos términos en cualquier organización. De la misma manera las tareas también se han descrito en inglés. 

Recordemos que Scrum no define como tal un método o herramienta para llevar el seguimiento del trabajo realizado. Scrum solo propone una serie de buenas prácticas por lo que el seguimiento se puede realizar con una simple hoja de excell, cuaderno o con herramientas de gestión especializadas. Yo he elegido el \textit{project board} de github dentro de mi repositorio del trabajo fin de máster.

Los roles o actores principales en Scrum serían los siguientes: 
\begin{itemize}	
	\item \textit{Scrum Master}: su trabajo prioritario es eliminar los obstáculos que impiden que el equipo no alcance el objetivo del \textit{sprint}. Esta persona no es el líder del equipo sino el nexo entre todos ellos. El autor del trabajo asume este rol.
	\item \textit{Product Owner}: representa a los \textit{stakeholders}\footnote{En toda organización, además de sus propietarios, participan diversos actores claves y grupos sociales que están constituidos por las personas o entes que, de una manera y otra, tienen interés en el desempeño de una empresa porque están relacionadas, bien directa, bien indirectamente, con ella. Fuente Wikipedia. \url{https://es.wikipedia.org/wiki/Parte_interesada_(empresas)}}, se asegura de que el equipo trabaje deforma adecuada desde la perspectiva de éstos. En un entorno real sería el encargado de escribir historias de usuario, priorizarlas y colocarla en el \textit{Product Backlog}. En este caso que se describe, recordemos en un entorno académico, el tutor académico sería el más indicado para asumir este rol. 
	\item \textit{Team} (equipo): en un entorno real sería el equipo que conforman todos los profesionales de diferentes ramas de conocimiento. El objetivo principal es el de entregar el producto. En este caso el autor asume el rol de equipo.
\end{itemize}

A continuación se describre el ciclo de desarrollo. Al inicio del proyecto el \textit{Product Owne}r debe definir los requisitos que serán los objetivos a cumplir. Éstos quedan reflejados y ordenados por orden de importancia en el \textit{Product Backlog}. Durante los sprints cada uno de estos objetivos serán subdivididos en tareas más pequeñas. En la figura \ref{backlog} se puede ver un ejemplo del \textit{Backlog}. 

\imagen{backlog}{\footnotesize{Detalle \textit{Product Backlog}.}\label{backlog}}

Al inicio de cada Sprint se extraen una serie de tareas de los elementos del \textit{Product Backlog} que conforman el \textit{Sprint Backlog}, las tareas son elegidas por el equipo. Estas tareas deben ser completadas durante el ciclo y no es posible añadir tareas nuevas a este \textit{Sprint} durante el tiempo que se está realizando el mismo.

\begin{itemize}	
	\item \textit{Product Backlog}: en un entorno real, consistiría en un listado de historias de usuario, obtenidas con el cliente, que se irán incorporando al producto a partir de incrementos sucesivos. En el proyecto descrito aquí éste se ve reflejado en las tareas (\textit{issues} en github) que se van creado a medida que son necesarias. Estás tareas se asignan después a cada sprint.
	\item \textit{Sprint}: una de las bases de los proyectos ágiles es el desarrollo mediante las iteraciones incrementales. En Scrum a cada iteración se le denomina \textit{Sprint}. En el caso concreto de este trabajo, excepto el último, todos tienen una duración de quince días.
\end{itemize}


A continuación se describen las reuniones que se realizan en el marco de desarrollo de trabajo \textit{Scrum}. 
\begin{itemize}	
	\item \textit{Daily Meeting}: se realiza una reunión breve diaria en la que cada miembro del equipo cuenta el estado en el que se encuenta la parte del proyecto que está realizando. En el caso del proyecto que se describe en esta documentación deberá ser una reflexión personal.
	\item \textit{Planning Meeting}: se realiza al inicio de cada \textit{Sprint} con el objetivo de escoger las tareas necesarias a realizar. Recodemos que se deben escoger del \textit{Product Backlog}.
	\item \textit{Review Meeting}:  esta reunión se realiza con el cliente al final de cada \textit{Sprint}, suele tener una duración máxima de cuatro horas y se revisa las historias de usuarios que han finalizado y las que no. Además se suele realizar una pequeña demo.
	\item \textit{Retrospective Meeting}: esta reunión también se realizala finalizar el \textit{Sprint}. Todos los miembros del equipo expresan sus opiniones sobre la iteracción finaliza. El objetivo de esta reunión es realizar una mejora continua del proceso. 
\end{itemize}


A continuación se describen los diferentes \emph{sprints} que se han realizado de manera detallada así como las planificaciones inicial y final.

%PlanificacionInicial-----------------------------
\subsubsection{Planificación inicial}

 \hl{Incluir aqui diagrama ganntttt INICIAL}

%\imagen{planificacioninicial}{\footnotesize{Planificación temporal inicial.}\label{planificacioninicial}

%Previo-----------------------------
\subsubsection{Mayo/Junio 2019}\label{previo}
El trabajo fin de máster fue elegido a finales de marzo de 2019. Debido a la intesa carga de trabajo del máster y a que la actividad del autor no se limita únicamente al ámbito académico sino que también desempeña un trabajo a tiempo completo durante la semana, no se pudo empezar a desarrollar una actividad completa hasta el mes de junio. 
En la figura \ref{scrumBoard} se puede ver el \textit{board} al inicio.

\imagen{scrumBoard}{\footnotesize{Detalle \textit{Scrum board}.}\label{scrumBoard}}

%Sprint0-----------------------------
\subsubsection{Sprint 0: 15/06/2019 - 01/07/2019}\label{sprint0}
\begin{itemize}
	\item[$\ast$] \underline{Duración}:  15 días
	\item[$\ast$] \underline{Descripción \textit{Sprint}}: Este primer \textit{Sprint} no se puede considerar como tal, de ahí el número cero, ya que no estaba al cien por cien en el trabajo fin de máster. Mi prioridad eran las asignaturas ya que sin ellas no podría desarrollarlo. Durante estos primeros quince días me dediqué a crear tareas y por tanto a llenar el \textit{backlog}. Si bien es cierto que la mayoría eran tareas como cosas de la documentación, cosas muy obvias sobre el desarrollo o tema académico. Durante la continuación del desarrollo se iba llenando el backlog con nuevas tareas que me iban surgiendo. De la misma manera también mucho tiempo dedicado a leer durante este \textit{Sprint}.
\end{itemize}
%Sprint1-----------------------------
\subsection{Sprint 1: 01/07/2019- 15/07/2019 (ver figura \ref{sprint1_img})}\label{sprint1}
\begin{itemize}
\item[$\ast$] \underline{Duración}:  15 días
\item[$\ast$] \underline{Descripción \textit{Sprint}}: En este primer \textit{Sprint} efectivo se comienza a trabajar en hacer un criba sobre todo lo investigado, además se comienza también a crear la documentación, repositorios y a hacer un análisis preliminar de los datos con los que vamos a trabajar.
\end{itemize}
\begin{enumerate}
	\item \textbf{\textit{Planning meeting}}
	\begin{description}
		\item[Reseña:] La primera reunión se fijan las tareas que se van a realizar este \textit{Sprint} y se concuerda con el tutor académico lo que tenemos en mente para comenzar con el desarrollo.
	\end{description}
	\item \textbf{\textit{Sprint planning}}
	\item[$-$] En esta reunión se establecen la tareas (\textit{user stories}) a desarrollar durante el \textit{Sprint}. 
	\begin{itemize}
		\item Crear repositorio.
		\item Investigar tecnologías principales.
		\item Analizar datos.
		\item Discutir propuesta con tutor académico.
		\item Discutir propuesta con tutor empresa.
		\item Terminar investigación sobre legislación.
		\item Terminar investigación científica sobre calidad del aire.
		\item Documentación: Introducción.
		\item Documentación: Objetivos.
	\end{itemize}
	\item \textbf{\textit{Retrospective meeting}}
	\begin{itemize}
		\item ¿Qué ha ido bien?
		\begin{itemize}
			\item Toda la parte de investigación tanto sobre legislación como científica.
			\item Reuniones muy productivas con los tutores para aclarar diversas dudas de índole diferente.
			\item Comienzo con la memoria de manera constante y paralela a las demás partes.
		\end{itemize}
		\item ¿Qué dificultades hemos encontrado?
		\begin{itemize}
			\item Debido a la enorme cantidad de información teórica que he recabado me va a ser complicado volcarla toda rápida
			\item Exiten numerosísimas opciones en cuanto a tecnologías. La oferta en el mercado es muy amplia y ello me crea  incertidumbre al respecto.
		\end{itemize}
	\end{itemize}	
\end{enumerate}
\imagen{sprint1}{\footnotesize{Detalle Sprint 1.}\label{sprint1_img}}

%Sprint2-----------------------------
\subsection{Sprint 2: 15/07/2019- 31/07/2019(ver figura \ref{sprint2_img})}\label{sprint2}
\begin{itemize}
\item[$\ast$] \underline{Duración}:  15 días
\item[$\ast$] \underline{Descripción \textit{Sprint}}: En este \textit{Sprint} se ha enfocado más a comprender, analizar e interpretar toda la información de la que se disponía. Ha sido un sprint de escribir mucho, asimilación y análisis antes de volver a la parte técnica. A veces es necesario tener claro todos los conceptos para saber como avanzar. 
\end{itemize}
\begin{enumerate}
\item \textbf{\textit{Planning meeting}}
\begin{description}
	\item[Reseña:] En esta reunión se fijan las tareas que se van a realizar este \textit{Sprint}. Además se añaden algunas nuevas al backlog como las referidas a la parte técnica que se realizarán en scripts sucesivos y que vienen derivadas de las necesidades de éste \textit{Sprint}.
\end{description}
\item \textbf{\textit{Sprint planning}}
	\item[$-$] En esta reunión se establecen la tareas (\textit{user stories}) a desarrollar durante el \textit{Sprint}. 
\begin{itemize}
	\item Creación de scripts y arquitectura del TFM.
	\item Ppreprocesamiento y limpieza de datos.
	\item Documentación: conceptos teóricos (Big Data).
	\item Documentación: conceptos teóricos (Legislación).
	\item Documentación: conceptos teóricos (Método análisis).
	\item Documentación: conceptos teóricos (Análisis científico).

\end{itemize}
\item \textbf{\textit{Retrospective meeting}}
\begin{itemize}
	\item ¿Qué ha ido bien?
	\begin{itemize}
		\item Con este sprint se ha comprendido una de las partes más importantes del trabajo: la parte científica.
		\item Toda la parte de preprocesamiento de datos.
	\end{itemize}
	\item ¿Qué dificultades hemos encontrado?
	\begin{itemize}
		\item Debido a la extensión de los conceptos teóricos no he podido progresar todo lo que me hubiera gustado en la parte técnica. Será necesario refactorizar código en \textit{sprints} sucesivos.
	\end{itemize}
\end{itemize}	
\end{enumerate}
\imagen{sprint2}{\footnotesize{Detalle Sprint 2.}\label{sprint2_img}}

%Sprint3-----------------------------
\subsection{Sprint 3: 31/07/2019 - 15/08/2019 (ver figura \ref{sprint3_img})}\label{sprint3}
\begin{itemize}
	\item[$\ast$] \underline{Duración}:  15 días
	\item[$\ast$] \underline{Descripción \textit{Sprint}}: Este tercer \textit{Sprint} se ha dirigido más a la parte técnica. Tanto seguir y terminar tareas del \textit{sprint }anterior como en la parte de visualización.
\end{itemize}
\begin{enumerate}
	\item \textbf{\textit{Planning meeting}}
	\begin{description}
		\item[Reseña:] En esta reunión he escogido las tareas a realizar en este sprint.
	\end{description}
	\item \textbf{\textit{Sprint planning}}
		\item[$-$] En esta reunión se establecen la tareas (\textit{user stories}) a desarrollar durante el \textit{Sprint}. 
	\begin{itemize}
		\item Visualización de datos en \textit{Python}.
		\item Visualización de datos en \textit{R}.
		\item Refactorización de código.
		\item Documentación: aspectos relevantes del desarrollo.
		\item Documentación: técnicas y herramientas.
		\item Documentación: trabajos relacionados.
		\item Enviar copia a tutores de los puntos de la memoria desarrollados en el \textit{sprint} anterior.
	\end{itemize}
	\item \textbf{\textit{Retrospective meeting}}
	\begin{itemize}
		\item ¿Qué ha ido bien?
		\begin{itemize}
			\item Este \textit{sprint} se ha enfocado en trabajar las visualizaciones tanto en \textit{Python} como en \textit{R}. 
			\item Se ha avanzado mucho en la memoria y además se ha enviado una primera copia a los tutores para su revisión.
		\end{itemize}
		\item ¿Qué dificultades hemos encontrado?
		\begin{itemize}
			\item Ninguna dificultad destacada. 
		\end{itemize}
	\end{itemize}	
\end{enumerate}
\imagen{sprint3}{\footnotesize{Detalle Sprint 3.}\label{sprint3_img}}

\subsection{Sprint 4: 15/08/2019 - 31/08/2019 (ver figura \ref{sprint4_img})}\label{sprint4}
\begin{itemize}
	\item[$\ast$] \underline{Duración}:  15 días
	\item[$\ast$] \underline{Descripción \textit{Sprint}}: Este cuarto \textit{Sprint} se ha dirigido más a la parte de la visualización final aunque sin olvidar la documentación.
\end{itemize}
\begin{enumerate}
	\item \textbf{\textit{Planning meeting}}
	\begin{description}
		\item[Reseña:] En esta reunión he escogido las tareas a realizar en este \textit{sprint}.
	\end{description}
	\item \textbf{\textit{Sprint planning}}
	\item[$-$] En esta reunión se establecen la tareas (\textit{user stories}) a desarrollar durante el \textit{Sprint}. 
	\begin{itemize}
		\item Visualización de datos con PowerBI.
		\item Investigación de uso de la herramienta software \textit{tinybird}.
		\item Discusión de resultados.
		\item Documentación: plan de proyecto software.
		\item Documentación: entorno experimental.
		\item Documentación: resultados.
		\item Enviar segunda copia de la memoria a los tutores
	\end{itemize}
	\item \textbf{\textit{Retrospective meeting}}
	\begin{itemize}
		\item ¿Qué ha ido bien?
		\begin{itemize}
			\item Este \textit{sprint} se ha dirigido a trabajar en la parte de la visualización final.
			\item Se ha avanzado en la memoria respecto de las partes finales de la misma. 
			\item Al introducirse una nueva herramienta (tinybird) se ha estado investigando su funcionamiento.
		\end{itemize}
		\item ¿Qué dificultades hemos encontrado?
		\begin{itemize}
			\item La parte de la visualización final se ha comido demasiado tiempo por lo que todas las tareas de documentación no han podido finalizarse tal y como se esperaba.
		\end{itemize}
	\end{itemize}	
\end{enumerate}
%\imagen{sprint4}{\footnotesize{Detalle Sprint 4.}\label{sprint4_img}}

\subsection{Sprint 5: 31/08/2019- 05/09/2019 (ver figura \ref{sprint5_img})}\label{sprint5}

%\imagen{sprint5}{\footnotesize{Detalle Sprint 5.}\label{sprint5_img}}

%PlanificacionFinal-----------------------------
\subsubsection{Planificación final}

\hl{Incluir aqui diagrama ganntttt FINAL}

%\imagen{planificacionfinal}{\footnotesize{Planificación temporal final.}\label{planificacionfinal}


\section{Viabilidad legal}\label{estudio-viabilidad}

La viabilidad legal se centra principalmente en el estudio de las licencias software utilizadas. Realizaremos una tabla resumen sobre las licencias. Indicar aquí que al tratarse de un estudio analítico se ha intentado seguir e utilizar software de carácter \textit{open source}.

\begin{table}[H]
	\begin{center}
		\begin{tabular}{|l|l|}
			\hline
			Dependencias & Licencia \\
			\hline \hline
			\emph{Github pages} & \href{https://opensource.org/licenses/MIT}{MIT}\\ \hline
			\emph{R} & \href{https://en.wikipedia.org/wiki/GNU_General_Public_License#Version_2}{GNU}\\ \hline
			\emph{Python} & \href{https://es.wikipedia.org/wiki/Python_Software_Foundation_License}{PSFL}\\ \hline
			\emph{Vega} & \href{https://github.com/vega/vega/blob/master/LICENSET}{BSD-3 Clause}\\ \hline
			
		\end{tabular}
		\caption{Tabla resumen-licencias.}
		\label{tabla:licencias}
	\end{center}
\end{table}



Una licencia software es un contrato entre el autor o titular de los derechos de explotación o distribución y el usuario consumidor, usuario profesional o empresa, para la utilización del software cumpliendo una serie de términos y condiciones establecidas dentro de sus cláusulas. Todo el software que empleo es su amplia mayoría es libre ya que la mayoría de código viene o bien de desarrolladores \emph{amateur} o es libre desde su creación. Emplean por tanto las licencias descritas en la tabla anterior por lo que diferencia entre ellas varían en términos como el nombramiento del autor o la garantía , son licencia comunes en software libre \cite{githublicense} , además github provee una página para ayudarte en la elección de tu licencia \citep{githubchoose}. La menos permisiva de entre las descritas es \emph{MIT} \citep{mit}  y la más es la licencia \emph{Apache-2.0} \cite{apache}.

\imagen{licencias}{\footnotesize{Gráfico licencias Open Source in GitHub. Fuente:  \url{https://cartograf.net}.}}

Para el proyecto de manera general (está alojado en github de forma libre) he elegido la licencia  \href{https://github.com/aguadotzn/TFM_MUINBDES/blob/master/LICENSE}{MIT}.

Para la documentación he elegido una licencia Creative commons, en concreto se ha elegido la \emph{Attribution-NonCommercial 4.0 International (CC BY-NC 4.0)}.

\imagen{creativecommons}{\footnotesize{Licencia Creative Commons.}}\label{commons}


\section{Links Importantes}\label{links_final}




\apendice{Links proyecto}

Este apéndice recoge  los links a los diferentes recursos del proyecto.


   \begin{table}[H]
	\begin{center}
		\begin{tabular}{|l|l|}
			\hline
			\textbf{Recurso} & \textit{Link} \\
			\hline \hline
			Repositorio github & \href{https://github.com/aguadotzn/TFM\_MUINBDES}{github/TFM\_MUINBDES}  \\ \hline
			Página web resumen & \href{https://aguadotzn.github.io/TFM_MUINBDES/}{Web}  \\ \hline
			Visualización interactiva & \href{https://app.powerbi.com/view?r=eyJrIjoiMmY1NGQ5YjYtYThiMS00YWJhLTllMDEtNmRhNTBhZWQzOTI1IiwidCI6IjJhYTNiMGI1LWE3ODItNGYzOC1hODk4LWU0ODNiMjBlOGQ2MSIsImMiOjh9}{PowerBI}  \\ \hline
			Mapa estaciones & \href{			https://aguadotzn.carto.com/builder/7a6bc6ca-594c-44ad-8bbe-add7757e0f0d/embed}{CartoDB Map}  \\ \hline
			


		\end{tabular}
		\caption{Enlaces a los diferentes recursos.}
	\end{center}
	\label{medidores}
\end{table}   




%%BibliografíaANEXOS
%%\bibliographystyle{plainnat}
%%\bibliography{bibliografiaAnexos}

\end{document}
